\documentclass[11pt, a4paper]{article}
\usepackage[english]{babel}
\usepackage[T1]{fontenc}
\usepackage{geometry}
\geometry{
    textwidth=6in,
    top=1in,
    bottom=1in,
    headheight=12pt,
    headsep=25pt,
    footskip=30pt
}
\usepackage{setspace}
\setstretch{1.2}
\usepackage{enumitem}
\usepackage{amsmath, amsthm, amssymb, amsfonts}
\usepackage{mathtools}
\usepackage{tikz-cd}
\usepackage[hidelinks]{hyperref}

\title{\textbf{$1^{\underline{st}}$ Part Problem Set for Applied Abstract Algebra} \\ Quivers: Representations and Morphisms}
\author{}
\date{}

\begin{document}
\maketitle

\begin{enumerate}[leftmargin=*, label=\textbf{\arabic*.}]

%=============================================================================
% Problem 1
%=============================================================================
\item Let $M, M' \in \operatorname{Rep} Q$. Show that the set of morphisms $\operatorname{Hom}(M, M')$ is a $k$-vector space.

\begin{proof}
Let $ M=(M_i, \phi_\alpha) $ and $ M'=(M_i', \psi_\alpha) $. $\forall f,g \in \operatorname{Hom}(M,M')$ since $ f,g $ are morphisms, for any arrow $ \alpha: i \to j $, we have $ \psi_\alpha f_i = f_j \phi_\alpha $ and $ \psi_\alpha g_i = g_j \phi_\alpha $.
    \begin{enumerate}
        \item \textbf{Addition:}
        \[
        \psi_\alpha (f+g)_i = \psi_\alpha (f_i + g_i) = \psi_\alpha f_i + \psi_\alpha g_i = f_j \phi_\alpha + g_j \phi_\alpha = (f_j + g_j) \phi_\alpha = (f+g)_j \phi_\alpha.
        \]
        Thus $ f+g \in \mathrm{Hom}(M,M') $.
        \item \textbf{Scalar Multiplication:}
        \[
        \psi_\alpha (\lambda f)_i = \psi_\alpha (\lambda f_i) = \lambda (\psi_\alpha f_i) = \lambda (f_j \phi_\alpha) = (\lambda f_j) \phi_\alpha = (\lambda f)_j \phi_\alpha.
        \]
        Thus $ \lambda f \in \mathrm{Hom}(M,M') $.
    \end{enumerate}
The zero morphism $ 0=(0)_{i \in Q_0} $ clearly satisfies the commutativity condition. The vector space axioms follow from the structure of linear maps between vector spaces.
\end{proof}

%=============================================================================
% Problem 2
%=============================================================================
\item Let $Q$ be the quiver
\[
1 \xrightarrow{\quad} 2 \xleftarrow{\quad} 3
\]
and consider the representations
\[
M\colon\quad k \xrightarrow{\begin{psmallmatrix} 1 \\ 0 \end{psmallmatrix}} k^2 \xleftarrow{\begin{psmallmatrix} 0 \\ 1 \end{psmallmatrix}} k
\qquad\qquad
M'\colon\quad k \xrightarrow{\begin{psmallmatrix} 1 \\ 0 \end{psmallmatrix}} k^2 \xleftarrow{\begin{psmallmatrix} 1 \\ 0 \end{psmallmatrix}} k
\]
\begin{enumerate}[label=(\alph*)]
    \item Show that $M$ and $M'$ are not indecomposable.
    \begin{proof}
    We exhibit a nontrivial direct sum decomposition for each representation.

    \medskip\noindent\textbf{Decomposition of $M$.}
    Define subrepresentations
    \[
    N_1\colon\quad k \xrightarrow{\begin{psmallmatrix} 1 \\ 0 \end{psmallmatrix}} \operatorname{span}(e_1) \xleftarrow{0} 0,
    \qquad
    N_2\colon\quad 0 \xrightarrow{0} \operatorname{span}(e_2) \xleftarrow{\begin{psmallmatrix} 0 \\ 1 \end{psmallmatrix}} k.
    \]
    At each vertex we have $M(1) = k = N_1(1) \oplus N_2(1)$, $M(2) = k^2 = \operatorname{span}(e_1) \oplus \operatorname{span}(e_2) = N_1(2) \oplus N_2(2)$, and $M(3) = k = N_1(3) \oplus N_2(3)$.
    The linear maps of $M$ respect this decomposition: the map $\binom{1}{0}\colon k \to k^2$ lands in $\operatorname{span}(e_1) = N_1(2)$, and $\binom{0}{1}\colon k \to k^2$ lands in $\operatorname{span}(e_2) = N_2(2)$.
    Hence $M = N_1 \oplus N_2 \cong P(1) \oplus P(3)$, so $M$ is not indecomposable.

    \medskip\noindent\textbf{Decomposition of $M'$.}
    Define subrepresentations
    \[
    N_1'\colon\quad k \xrightarrow{1} \operatorname{span}(e_1) \xleftarrow{1} k,
    \qquad
    N_2'\colon\quad 0 \xrightarrow{0} \operatorname{span}(e_2) \xleftarrow{0} 0.
    \]
    At each vertex: $M'(1) = k = N_1'(1) \oplus N_2'(1)$, $M'(2) = k^2 = \operatorname{span}(e_1) \oplus \operatorname{span}(e_2) = N_1'(2) \oplus N_2'(2)$, and $M'(3) = k = N_1'(3) \oplus N_2'(3)$.
    Both arrows $\binom{1}{0}\colon M'(1) \to M'(2)$ and $\binom{1}{0}\colon M'(3) \to M'(2)$ land in $\operatorname{span}(e_1) = N_1'(2)$, so the maps respect the decomposition.
    Hence $M' = N_1' \oplus N_2' \cong I(2) \oplus S(2)$, so $M'$ is not indecomposable.
    \end{proof}
    \item Show that $M$ and $M'$ are not isomorphic.
    \begin{proof}
    Suppose for contradiction that $f\colon M \to M'$ is an isomorphism, so $f = (f_1, f_2, f_3)$ with $f_i$ invertible. The commutativity conditions give:
    \[
    f_2 \begin{psmallmatrix} 1 \\ 0 \end{psmallmatrix} = \begin{psmallmatrix} 1 \\ 0 \end{psmallmatrix} f_1,
    \qquad
    f_2 \begin{psmallmatrix} 0 \\ 1 \end{psmallmatrix} = \begin{psmallmatrix} 1 \\ 0 \end{psmallmatrix} f_3.
    \]
    Let $f_1 = (a) \in k^*$, $f_3 = (c) \in k^*$, and $f_2 = \begin{psmallmatrix} p & q \\ r & s \end{psmallmatrix} \in \mathrm{GL}_2(k)$.

    From the first equation: $f_2 \binom{1}{0} = \binom{p}{r} = \binom{a}{0}$, so $p = a$ and $r = 0$.

    From the second equation: $f_2 \binom{0}{1} = \binom{q}{s} = \binom{c}{0}$, so $q = c$ and $s = 0$.

    Then $f_2 = \begin{psmallmatrix} a & c \\ 0 & 0 \end{psmallmatrix}$, which has $\det(f_2) = 0$. This contradicts $f_2 \in \mathrm{GL}_2(k)$.

    Therefore $M \not\cong M'$.
    \end{proof}
\end{enumerate}

%=============================================================================
% Problem 3
%=============================================================================
\item Let $Q$ be the quiver
\[
1 \xrightarrow{\quad} 2 \xleftarrow{\quad} 3
\]
and let $M$ be the representation
\[
M\colon\quad k^2 \xrightarrow{\begin{psmallmatrix} 1 & 0 \\ 0 & 1 \\ 0 & 0 \end{psmallmatrix}} k^3 \xleftarrow{\begin{psmallmatrix} 1 \\ 0 \\ 0 \end{psmallmatrix}} k\,.
\]
\begin{enumerate}[label=(\alph*)]
    \item Write $M$ as a direct sum of the indecomposable representations listed below:
    \begin{align*}
        S(1) &\colon\ k \to 0 \leftarrow 0 &
        S(2) &\colon\ 0 \to k \leftarrow 0 &
        S(3) &\colon\ 0 \to 0 \leftarrow k \\[6pt]
        P(1) &\colon\ k \xrightarrow{1} k \leftarrow 0 &
        I(2) &\colon\ k \xrightarrow{1} k \xleftarrow{1} k &
        P(3) &\colon\ 0 \to k \xleftarrow{1} k
    \end{align*}
    \begin{proof}
    We claim that
    \[
    M \cong I(2) \oplus P(1) \oplus S(2).
    \]
    
    The dimension vectors of the summands are
    \[
    \underline{\dim}\,I(2) = (1,1,1),\qquad \underline{\dim}\,P(1) = (1,1,0),\qquad \underline{\dim}\,S(2) = (0,1,0),
    \]
    and their sum is
    \[
    (1,1,1)+(1,1,0)+(0,1,0) = (2,3,1) = \underline{\dim}\,M.
    \]
    
    We now verify that the maps of $I(2)\oplus P(1)\oplus S(2)$ agree with those of $M$. Order the components at each vertex as:
    \[
    \text{Vertex }1\colon\; I(2)(1)\oplus P(1)(1) = k\oplus k = k^2,
    \]
    \[
    \text{Vertex }2\colon\; I(2)(2)\oplus P(1)(2)\oplus S(2)(2) = k\oplus k\oplus k = k^3,
    \]
    \[
    \text{Vertex }3\colon\; I(2)(3) = k.
    \]
    
    \medskip\noindent\textbf{Map at arrow $1\to 2$.}
    In the direct sum, the map $k^2 \to k^3$ sends
    \[
    (a,b) \longmapsto \bigl(I(2)(\alpha)(a),\; P(1)(\alpha)(b),\; 0\bigr) = (a,\,b,\,0),
    \]
    since both $I(2)$ and $P(1)$ have identity maps at this arrow, and $S(2)(1)=0$ contributes nothing. In matrix form this gives
    \[
    \begin{pmatrix} 1 & 0 \\ 0 & 1 \\ 0 & 0 \end{pmatrix},
    \]
    which agrees with the map of $M$.
    
    \medskip\noindent\textbf{Map at arrow $3\to 2$.}
    In the direct sum, the map $k \to k^3$ sends
    \[
    c \longmapsto \bigl(I(2)(\beta)(c),\; 0,\; 0\bigr) = (c,\,0,\,0),
    \]
    since $I(2)$ has identity map at this arrow, and $P(1)(3)=S(2)(3)=0$ contribute nothing. In matrix form this gives
    \[
    \begin{pmatrix} 1 \\ 0 \\ 0 \end{pmatrix},
    \]
    which agrees with the map of $M$. Therefore $M \cong I(2) \oplus P(1) \oplus S(2)$.
    \end{proof}
    \item Show that there is a non-split short exact sequence
    \[
        0 \longrightarrow X \longrightarrow Y \longrightarrow Z \longrightarrow 0\,.
    \]
    \begin{proof}
    Consider the representations
    \[
    S(2) = (0,\,k,\,0),\qquad P(1) = (k,\,k,\,0)\ \text{with map } \mathrm{id}_k,\qquad S(1) = (k,\,0,\,0).
    \]
    
    Define morphisms $\iota\colon S(2)\to P(1)$ and $p\colon P(1)\to S(1)$ by
    \[
    \iota = (\iota_1,\,\iota_2,\,\iota_3) = (0,\;\mathrm{id}_k,\;0),\qquad p = (p_1,\,p_2,\,p_3) = (\mathrm{id}_k,\;0,\;0).
    \]
    
    \medskip\noindent\textbf{Morphism verification.}
    The only nontrivial arrow is $\alpha\colon 1\to 2$ with $P(1)(\alpha)=\mathrm{id}_k$. For~$\iota$, the commutativity diagram at $\alpha$ is
    \[
    \begin{tikzcd}
    S(2)(1) = 0 \arrow[r, "0"] \arrow[d, "\iota_1=0"'] & S(2)(2) = k \arrow[d, "\iota_2=\mathrm{id}_k"] \\
    P(1)(1) = k \arrow[r, "\mathrm{id}_k"'] & P(1)(2) = k
    \end{tikzcd}
    \]
    Both compositions give $0$, so the diagram commutes. For~$p$:
    \[
    \begin{tikzcd}
    P(1)(1) = k \arrow[r, "\mathrm{id}_k"] \arrow[d, "p_1=\mathrm{id}_k"'] & P(1)(2) = k \arrow[d, "p_2=0"] \\
    S(1)(1) = k \arrow[r, "0"'] & S(1)(2) = 0
    \end{tikzcd}
    \]
    Both compositions give $0$, so the diagram commutes.
    
    \medskip\noindent\textbf{Exactness.}
    
    $\iota$ is injective: the only nontrivial component is $\iota_2=\mathrm{id}_k$, which is injective.
    
    $p$ is surjective: $p_1=\mathrm{id}_k$ is surjective, and $p_2\colon k\to 0$, $p_3\colon 0\to 0$ are trivially surjective.
    
    Exactness at $P(1)$: we verify $\ker p = \operatorname{im}\iota$ at each vertex.
    \begin{itemize}
    \item Vertex~$1$: $\ker(\mathrm{id}_k)=0=\operatorname{im}(0)$.
    \item Vertex~$2$: $\ker(0\colon k\to 0)=k=\operatorname{im}(\mathrm{id}_k)$.
    \item Vertex~$3$: both are $0$.
    \end{itemize}
    
    Hence
    \[
    0 \longrightarrow S(2) \xrightarrow{\;\iota\;} P(1) \xrightarrow{\;p\;} S(1) \longrightarrow 0
    \]
    is a short exact sequence.
    
    \medskip\noindent\textbf{Non-split.}
    If the sequence splits, then
    \[
    P(1) \cong S(1) \oplus S(2).
    \]
    However, the direct sum $S(1)\oplus S(2) = (k,\,k,\,0)$ has map $0\oplus 0 = 0$ at arrow $1\to 2$, while $P(1)$ has map $\mathrm{id}_k\neq 0$. This is a contradiction, so the sequence does not split.
    \end{proof}
\end{enumerate}

%=============================================================================
% Problem 4
%=============================================================================
\item Find all indecomposable representations up to isomorphism of the quiver $1 \rightarrow 2$.

\begin{proof}
The indecomposable representations of $Q\colon 1\to 2$ are, up to isomorphism:
\[
S(1)\colon\ k \xrightarrow{0} 0,\qquad S(2)\colon\ 0 \xrightarrow{0} k,\qquad P(1)\colon\ k \xrightarrow{\mathrm{id}} k.
\]

\medskip\noindent\textbf{Step 1: Indecomposability.}

$S(1)$ and $S(2)$ have dimension vectors $(1,0)$ and $(0,1)$ respectively. Since no component can be split further, these are clearly indecomposable.

For $P(1)$ with dimension vector $(1,1)$, suppose
\[
P(1) = N_1 \oplus N_2
\]
is a nontrivial decomposition. Then we must have
\[
\underline{\dim}\,N_1 = (1,0),\qquad \underline{\dim}\,N_2 = (0,1)
\]
(or vice versa). In this decomposition, the map $\phi\colon P(1)(1)=k \to P(1)(2)=k$ would decompose as
\[
\phi = \phi|_{N_1} \oplus \phi|_{N_2} = 0 \oplus 0 = 0,
\]
since $\phi|_{N_1}\colon k \to 0$ and $\phi|_{N_2}\colon 0 \to k$ are both the zero map. This contradicts $\phi=\mathrm{id}_k\neq 0$. Hence $P(1)$ is indecomposable.

\medskip\noindent\textbf{Step 2: Completeness.}

Let
\[
M = (V_1,\,V_2,\,\phi)
\]
be any representation with $\phi\colon V_1\to V_2$. Set
\[
n_1 = \dim V_1,\qquad n_2 = \dim V_2,\qquad r = \operatorname{rank}(\phi).
\]

Choose a basis $\{e_1,\ldots,e_{n_1}\}$ of $V_1$ such that $\{e_1,\ldots,e_r\}$ maps to linearly independent vectors in $V_2$ and $\{e_{r+1},\ldots,e_{n_1}\}$ spans $\ker\phi$. Let
\[
f_i = \phi(e_i) \quad\text{for } 1\leq i\leq r,
\]
and extend $\{f_1,\ldots,f_r\}$ to a basis $\{f_1,\ldots,f_{n_2}\}$ of $V_2$. In these bases,
\[
\phi = \begin{pmatrix} I_r & 0 \\ 0 & 0 \end{pmatrix}.
\]

Then $M$ decomposes as:
\begin{itemize}
\item $r$ copies of $P(1)$: the $i$-th copy is spanned by $e_i$ at vertex~$1$ and $f_i$ at vertex~$2$, with $\phi(e_i)=f_i$, for $1\leq i\leq r$.
\item $(n_1-r)$ copies of $S(1)$: spanned by $e_{r+j}$ at vertex~$1$, with $\phi(e_{r+j})=0$.
\item $(n_2-r)$ copies of $S(2)$: spanned by $f_{r+l}$ at vertex~$2$.
\end{itemize}

Hence
\[
M \cong P(1)^{\oplus r} \oplus S(1)^{\oplus(n_1-r)} \oplus S(2)^{\oplus(n_2-r)}.
\]
Since every representation decomposes into copies of $S(1)$, $S(2)$, and $P(1)$, these are all the indecomposable representations up to isomorphism.
\end{proof}

%=============================================================================
% Problem 5
%=============================================================================
\item Let $Q$ be the quiver
\[
    \begin{tikzcd}
        1 \arrow[rd, "\alpha"] &   &                       \\
                               & 3 & 4 \arrow[l, "\gamma"] \\
        2 \arrow[ru, "\beta"]  &   &                      
        \end{tikzcd}
\]
and consider the following representations
\[
M\colon\quad
\begin{tikzcd}
    k \arrow[rd, "\begin{bmatrix} 1 \\ 0 \end{bmatrix}"] &     &                                                     \\
                                                         & k^2 & k \arrow[l, "\begin{bmatrix} 1 \\ 1 \end{bmatrix}"] \\
    k \arrow[ru, "\begin{bmatrix} 0 \\ 1 \end{bmatrix}"] &     &                                                    
    \end{tikzcd}
\]
\[
M'\colon\quad
\begin{tikzcd}
    0 \arrow[rd] &   &             \\
                 & k & 0 \arrow[l] \\
    0 \arrow[ru] &   &            
    \end{tikzcd}
\]  
\[
M''\colon\quad
\begin{tikzcd}
    k \arrow[rd, "1"] &   &                  \\
                      & k & k \arrow[l, "1"] \\
    k \arrow[ru, "1"] &   &                 
    \end{tikzcd}
\]
\begin{enumerate}[label=(\alph*)]
    \item Show that
    \[
        \operatorname{Hom}(M, M') = 0, \qquad \operatorname{Hom}(M', M) \cong k^2
    \]
    \[
        \operatorname{Hom}(M, M'') \cong k^2, \qquad \operatorname{Hom}(M'', M) = 0.
    \]
    \begin{proof}
    The quiver $Q$ has arrows $\alpha\colon 1\to 3$, $\beta\colon 2\to 3$, $\gamma\colon 4\to 3$. A morphism $f\colon A\to B$ consists of linear maps $(f_1,f_2,f_3,f_4)$ satisfying
    \[
    f_3\circ A(\delta) = B(\delta)\circ f_s
    \]
    for each arrow $\delta\colon s\to 3$.

    \medskip\noindent$\boldsymbol{\operatorname{Hom}(M,M')=0.}$
    
    A morphism $f\colon M\to M'$ has components
    \[
    f_1\colon k\to 0,\qquad f_2\colon k\to 0,\qquad f_3\colon k^2\to k,\qquad f_4\colon k\to 0.
    \]
    Since the codomain is $0$, we have $f_1=f_2=f_4=0$. The commutativity condition at $\alpha$ gives
    \[
    \begin{tikzcd}
    M(1) = k \arrow[r, "\binom{1}{0}"] \arrow[d, "f_1=0"'] & M(3) = k^2 \arrow[d, "f_3"] \\
    M'(1) = 0 \arrow[r, "0"'] & M'(3) = k
    \end{tikzcd}
    \qquad\Longrightarrow\qquad
    f_3\begin{pmatrix}1\\0\end{pmatrix} = 0.
    \]
    Similarly, the commutativity condition at $\beta$ gives
    \[
    f_3\begin{pmatrix}0\\1\end{pmatrix} = 0.
    \]
    Since $\binom{1}{0}$ and $\binom{0}{1}$ form a basis of $k^2$, we conclude $f_3=0$. Therefore $f=0$.

    \medskip\noindent$\boldsymbol{\operatorname{Hom}(M',M)\cong k^2.}$
    
    A morphism $f\colon M'\to M$ has components
    \[
    f_1\colon 0\to k,\qquad f_2\colon 0\to k,\qquad f_3\colon k\to k^2,\qquad f_4\colon 0\to k.
    \]
    Since the domains of $f_1,f_2,f_4$ are $0$, these maps are all zero. For each arrow $\delta\colon s\to 3$, the commutativity condition becomes
    \[
    f_3 \circ M'(\delta) = M(\delta)\circ f_s.
    \]
    Since $M'(s)=0$ for $s=1,2,4$, both $M'(\delta)$ and $f_s$ are the zero map, so the condition is automatically satisfied. Hence $f$ is completely determined by the arbitrary linear map
    \[
    f_3 \in \operatorname{Hom}_k(k,\,k^2) \cong k^2.
    \]

    \medskip\noindent$\boldsymbol{\operatorname{Hom}(M,M'')\cong k^2.}$
    
    A morphism $f\colon M\to M''$ has components
    \[
    f_1 = a\colon k\to k,\quad f_2 = b\colon k\to k,\quad f_3 = \begin{pmatrix}c&d\end{pmatrix}\colon k^2\to k,\quad f_4 = e\colon k\to k,
    \]
    where $a,b,c,d,e\in k$. The commutativity conditions give:
    
    At $\alpha\colon 1\to 3$:
    \[
    f_3\circ M(\alpha) = M''(\alpha)\circ f_1 \quad\Longrightarrow\quad \begin{pmatrix}c&d\end{pmatrix}\begin{pmatrix}1\\0\end{pmatrix} = 1\cdot a \quad\Longrightarrow\quad c = a.
    \]
    
    At $\beta\colon 2\to 3$:
    \[
    f_3\circ M(\beta) = M''(\beta)\circ f_2 \quad\Longrightarrow\quad \begin{pmatrix}c&d\end{pmatrix}\begin{pmatrix}0\\1\end{pmatrix} = 1\cdot b \quad\Longrightarrow\quad d = b.
    \]
    
    At $\gamma\colon 4\to 3$:
    \[
    f_3\circ M(\gamma) = M''(\gamma)\circ f_4 \quad\Longrightarrow\quad \begin{pmatrix}c&d\end{pmatrix}\begin{pmatrix}1\\1\end{pmatrix} = 1\cdot e \quad\Longrightarrow\quad e = a+b.
    \]
    
    Hence every morphism $f$ is determined by the pair $(a,b)\in k^2$, with $c=a$, $d=b$, $e=a+b$. Therefore
    \[
    \operatorname{Hom}(M,M'') \cong k^2.
    \]

    \medskip\noindent$\boldsymbol{\operatorname{Hom}(M'',M)=0.}$
    
    A morphism $f\colon M''\to M$ has components
    \[
    f_1 = a\colon k\to k,\quad f_2 = b\colon k\to k,\quad f_3 = \begin{pmatrix}c\\d\end{pmatrix}\colon k\to k^2,\quad f_4 = e\colon k\to k.
    \]
    
    At $\alpha\colon 1\to 3$:
    \[
    M(\alpha)\circ f_1 = f_3\circ M''(\alpha) \quad\Longrightarrow\quad \begin{pmatrix}1\\0\end{pmatrix}a = \begin{pmatrix}c\\d\end{pmatrix}\cdot 1 \quad\Longrightarrow\quad c = a,\; d = 0.
    \]
    
    At $\beta\colon 2\to 3$:
    \[
    M(\beta)\circ f_2 = f_3\circ M''(\beta) \quad\Longrightarrow\quad \begin{pmatrix}0\\1\end{pmatrix}b = \begin{pmatrix}c\\d\end{pmatrix} = \begin{pmatrix}a\\0\end{pmatrix} \quad\Longrightarrow\quad a = 0,\; b = 0.
    \]
    
    At $\gamma\colon 4\to 3$:
    \[
    M(\gamma)\circ f_4 = f_3\circ M''(\gamma) \quad\Longrightarrow\quad \begin{pmatrix}1\\1\end{pmatrix}e = \begin{pmatrix}a\\0\end{pmatrix} = \begin{pmatrix}0\\0\end{pmatrix} \quad\Longrightarrow\quad e = 0.
    \]
    
    All components vanish, so $\operatorname{Hom}(M'',M) = 0$.
    \end{proof}
    \item Show that $M$ is not isomorphic to $M' \oplus M''$.
    \begin{proof}
    Suppose for contradiction that
    \[
    M \cong M' \oplus M''.
    \]
    Then by the additivity of $\operatorname{Hom}$,
    \[
    \operatorname{Hom}(M,\,M') \cong \operatorname{Hom}(M'\oplus M'',\, M') \cong \operatorname{Hom}(M',\,M') \oplus \operatorname{Hom}(M'',\,M').
    \]
    Since $M'\neq 0$, the identity morphism $\mathrm{id}_{M'}$ is a nonzero element of $\operatorname{Hom}(M',M')$. Hence
    \[
    \operatorname{Hom}(M',\,M') \neq 0,
    \]
    which implies
    \[
    \operatorname{Hom}(M,\,M') \neq 0.
    \]
    This contradicts part~(a), where we showed $\operatorname{Hom}(M,M')=0$. Therefore
    \[
    M \not\cong M' \oplus M''.
    \]
    \end{proof}
    \item Show that there is a short exact sequence:
    \[
        0 \longrightarrow M' \longrightarrow M \longrightarrow M'' \longrightarrow 0\,.
    \]
    \begin{proof}
    We construct explicit morphisms $\iota\colon M'\to M$ and $p\colon M\to M''$ such that
    \[
    0 \longrightarrow M' \xrightarrow{\;\iota\;} M \xrightarrow{\;p\;} M'' \longrightarrow 0
    \]
    is exact.
    
    \medskip\noindent\textbf{Definition of $\iota\colon M'\to M$.}
    Since $M'(1)=M'(2)=M'(4)=0$, we set
    \[
    \iota_1 = 0\colon 0\to k,\qquad \iota_2 = 0\colon 0\to k,\qquad \iota_4 = 0\colon 0\to k,
    \]
    and define
    \[
    \iota_3 = \begin{pmatrix}1\\-1\end{pmatrix}\colon k\to k^2.
    \]
    Since all maps of $M'$ are zero, for each arrow $\delta\colon s\to 3$, the commutativity condition
    \[
    \iota_3 \circ M'(\delta) = M(\delta)\circ \iota_s
    \]
    becomes $0 = 0$, which is automatic. Moreover, $\iota_3$ is injective (since $\binom{1}{-1}\neq 0$), so $\iota$ is injective.
    
    \medskip\noindent\textbf{Definition of $p\colon M\to M''$.}
    Define
    \[
    p_1 = 1\colon k\to k,\qquad p_2 = 1\colon k\to k,\qquad p_3 = \begin{pmatrix}1&1\end{pmatrix}\colon k^2\to k,\qquad p_4 = 2\colon k\to k.
    \]
    We verify commutativity at each arrow.
    
    At $\alpha\colon 1\to 3$:
    \[
    p_3\circ M(\alpha) = \begin{pmatrix}1&1\end{pmatrix}\begin{pmatrix}1\\0\end{pmatrix} = 1 = 1\cdot 1 = M''(\alpha)\circ p_1.
    \]
    
    At $\beta\colon 2\to 3$:
    \[
    p_3\circ M(\beta) = \begin{pmatrix}1&1\end{pmatrix}\begin{pmatrix}0\\1\end{pmatrix} = 1 = 1\cdot 1 = M''(\beta)\circ p_2.
    \]
    
    At $\gamma\colon 4\to 3$:
    \[
    p_3\circ M(\gamma) = \begin{pmatrix}1&1\end{pmatrix}\begin{pmatrix}1\\1\end{pmatrix} = 2 = 1\cdot 2 = M''(\gamma)\circ p_4.
    \]
    Each $p_i$ is a nonzero scalar, so $p$ is surjective at every vertex.
    
    \medskip\noindent\textbf{Exactness at $M$.}
    
    First, $p\circ\iota = 0$: at vertex~$3$,
    \[
    p_3\circ\iota_3 = \begin{pmatrix}1&1\end{pmatrix}\begin{pmatrix}1\\-1\end{pmatrix} = 1+(-1) = 0.
    \]
    At other vertices, $\iota_i = 0$, so $(p\circ\iota)_i = p_i\circ 0 = 0$.
    
    Next, we verify $\ker p = \operatorname{im}\iota$ at each vertex:
    \begin{itemize}
    \item Vertices $1,2,4$: each $p_i$ is a nonzero scalar (hence injective), so $\ker p_i = 0$. Also $\iota_i = 0$, so $\operatorname{im}\iota_i = 0$.
    \item Vertex~$3$:
    \[
    \ker p_3 = \left\{(x,y)\in k^2 : x+y = 0\right\} = \left\langle\begin{pmatrix}1\\-1\end{pmatrix}\right\rangle = \operatorname{im}\iota_3.
    \]
    \end{itemize}
    
    Therefore
    \[
    0 \longrightarrow M' \xrightarrow{\;\iota\;} M \xrightarrow{\;p\;} M'' \longrightarrow 0
    \]
    is a short exact sequence.
    \end{proof}
\end{enumerate}

%=============================================================================
% Problem 6
%=============================================================================
\item Let $Q$ be the quiver
\[
\begin{tikzcd}
1 & 2 \arrow[l, "\alpha"', shift right] \arrow[l, "\beta", shift left]
\end{tikzcd}
\]
For any $\lambda \in k \cup \{\infty\}$, define $M_\lambda$ to be the representation:
\[
M_\lambda\colon\quad
\begin{tikzcd}
k & k \arrow[l, "\lambda"', shift right] \arrow[l, "1", shift left]
\end{tikzcd}
\quad \text{if } \lambda \in k;
\qquad\qquad
M_\infty\colon\quad
\begin{tikzcd}
k & k \arrow[l, "1"', shift right] \arrow[l, "0", shift left]
\end{tikzcd}
\quad \text{if } \lambda = \infty.
\]
\begin{enumerate}[label=(\alph*)]
    \item Show that each $M_\lambda$ is indecomposable.
    \begin{proof}
    Each $M_\lambda$ has dimension vector
    \[
    \underline{\dim}\,M_\lambda = (1,1).
    \]
    Suppose for contradiction that
    \[
    M_\lambda = N_1 \oplus N_2
    \]
    is a nontrivial decomposition (both $N_1,N_2\neq 0$). Since $\underline{\dim}\,N_1 + \underline{\dim}\,N_2 = (1,1)$, the only possibility is
    \[
    \underline{\dim}\,N_1 = (1,0),\qquad \underline{\dim}\,N_2 = (0,1)
    \]
    (or vice versa). The summand $N_2$ with $N_2(1)=0$ and $N_2(2)=k$ must be a subrepresentation. Since both arrows go from vertex~$2$ to vertex~$1$, this requires
    \[
    M_\lambda(\alpha)(k) \subseteq N_2(1) = 0 \qquad\text{and}\qquad M_\lambda(\beta)(k) \subseteq N_2(1) = 0.
    \]

    For $\lambda\in k$: the map $M_\lambda(\beta)=1\colon k\to k$ is nonzero, so
    \[
    M_\lambda(\beta)(k) = k \neq 0.
    \]
    This is a contradiction.

    For $\lambda=\infty$: the map $M_\infty(\alpha)=1\colon k\to k$ is nonzero, so
    \[
    M_\infty(\alpha)(k) = k \neq 0.
    \]
    This is again a contradiction.

    In both cases, no such subrepresentation $N_2$ exists, so $M_\lambda$ is indecomposable.
    \end{proof}
    \item Show that $M_\lambda \cong M_\mu$ if and only if $\lambda = \mu$. In particular, the number of indecomposable representations depends on the choice of the field $k$.
    \begin{proof}
    ($\Leftarrow$) If $\lambda=\mu$, the identity maps $f_1=f_2=\mathrm{id}_k$ give an isomorphism.

    \medskip
    ($\Rightarrow$) Let $f=(f_1,f_2)\colon M_\lambda\to M_\mu$ be an isomorphism, with
    \[
    f_1 = a\in k^*,\qquad f_2 = b\in k^*.
    \]
    The commutativity conditions for the two arrows $\alpha,\beta\colon 2\to 1$ are
    \[
    f_1\circ M_\lambda(\alpha) = M_\mu(\alpha)\circ f_2,\qquad f_1\circ M_\lambda(\beta) = M_\mu(\beta)\circ f_2.
    \]

    \textbf{Case 1:} $\lambda,\mu\in k$. Then $M_\lambda(\alpha)=\lambda$, $M_\lambda(\beta)=1$, $M_\mu(\alpha)=\mu$, $M_\mu(\beta)=1$. The conditions become
    \[
    a\lambda = \mu b,\qquad a\cdot 1 = 1\cdot b.
    \]
    From the second equation, $a=b$. Substituting into the first:
    \[
    a\lambda = \mu a.
    \]
    Since $a\neq 0$, we can divide both sides by $a$ to get $\lambda=\mu$.

    \textbf{Case 2:} $\lambda\in k$, $\mu=\infty$. Then $M_\lambda(\beta)=1$ and $M_\infty(\beta)=0$. The commutativity at $\beta$ gives
    \[
    a\cdot 1 = 0\cdot b = 0,
    \]
    so $a=0$, contradicting $a\in k^*$. Hence no isomorphism exists.

    The case $\lambda=\infty$, $\mu\in k$ is symmetric.

    \textbf{Case 3:} $\lambda=\mu=\infty$. Trivially $\lambda=\mu$.

    \medskip
    In all cases, $M_\lambda\cong M_\mu$ implies $\lambda=\mu$. Since
    \[
    |k\cup\{\infty\}| = |k|+1,
    \]
    the number of pairwise non-isomorphic indecomposable representations of this form depends on the cardinality of the field $k$.
    \end{proof}
    \item Show that $\operatorname{Hom}(M_\lambda, M_\mu) = 0$ if $\lambda \neq \mu$.
    \begin{proof}
    Let $f=(f_1,f_2)\colon M_\lambda\to M_\mu$ be a morphism, with $f_1=a$ and $f_2=b$ (scalars). The commutativity conditions are
    \[
    f_1\circ M_\lambda(\alpha) = M_\mu(\alpha)\circ f_2,\qquad f_1\circ M_\lambda(\beta) = M_\mu(\beta)\circ f_2.
    \]

    \textbf{Case 1:} $\lambda,\mu\in k$. The conditions become
    \[
    a\lambda = \mu b,\qquad a = b.
    \]
    Substituting $b=a$ into the first equation:
    \[
    a\lambda = \mu a \quad\Longrightarrow\quad a(\lambda-\mu) = 0.
    \]
    Since $\lambda\neq\mu$, we have $\lambda-\mu\neq 0$, so $a=0$. Then $b=a=0$.

    \textbf{Case 2:} $\lambda\in k$, $\mu=\infty$. Then $M_\lambda(\beta)=1$, $M_\infty(\beta)=0$. The condition at $\beta$ gives
    \[
    a\cdot 1 = 0\cdot b = 0,
    \]
    so $a=0$. The condition at $\alpha$ then gives
    \[
    0\cdot\lambda = 1\cdot b \quad\Longrightarrow\quad b = 0.
    \]

    \textbf{Case 3:} $\lambda=\infty$, $\mu\in k$. Then $M_\infty(\beta)=0$, $M_\mu(\beta)=1$. The condition at $\beta$ gives
    \[
    a\cdot 0 = 1\cdot b \quad\Longrightarrow\quad b = 0.
    \]
    The condition at $\alpha$ gives
    \[
    a\cdot 1 = \mu\cdot 0 = 0 \quad\Longrightarrow\quad a = 0.
    \]

    In every case, $f_1=f_2=0$, so $f=0$. Therefore $\operatorname{Hom}(M_\lambda,M_\mu) = 0$.
    \end{proof}
    \item Show that for each $\lambda$ there is a short exact sequence
    \[
        0 \longrightarrow S(1) \longrightarrow M_\lambda \longrightarrow S(2) \longrightarrow 0\,,
    \]
    where $S(1)$ and $S(2)$ are the representations
    \[
        S(1)\colon\quad
        \begin{tikzcd}
        k & 0 \arrow[l, shift right] \arrow[l, shift left]
        \end{tikzcd}
        \qquad\qquad
        S(2)\colon\quad
        \begin{tikzcd}
        0 & k \arrow[l, shift right] \arrow[l, shift left]
        \end{tikzcd}
    \]
    \begin{proof}
    Define $\iota\colon S(1)\to M_\lambda$ and $p\colon M_\lambda\to S(2)$ by
    \[
    \iota_1 = \mathrm{id}_k\colon k\to k,\qquad \iota_2 = 0\colon 0\to k,
    \]
    \[
    p_1 = 0\colon k\to 0,\qquad p_2 = \mathrm{id}_k\colon k\to k.
    \]
    
    \medskip\noindent\textbf{Morphism verification.}
    Both arrows $\alpha,\beta$ go from vertex~$2$ to vertex~$1$. For $\iota$, the commutativity diagram at each arrow $\delta\in\{\alpha,\beta\}$ is
    \[
    \begin{tikzcd}
    S(1)(2) = 0 \arrow[r, "0"] \arrow[d, "\iota_2=0"'] & S(1)(1) = k \arrow[d, "\iota_1=\mathrm{id}_k"] \\
    M_\lambda(2) = k \arrow[r, "M_\lambda(\delta)"'] & M_\lambda(1) = k
    \end{tikzcd}
    \]
    Both compositions give $0$, so the diagram commutes.
    
    For $p$, the commutativity diagram at each arrow $\delta$ is
    \[
    \begin{tikzcd}
    M_\lambda(2) = k \arrow[r, "M_\lambda(\delta)"] \arrow[d, "p_2=\mathrm{id}_k"'] & M_\lambda(1) = k \arrow[d, "p_1=0"] \\
    S(2)(2) = k \arrow[r, "0"'] & S(2)(1) = 0
    \end{tikzcd}
    \]
    Both compositions give $0$, so the diagram commutes.
    
    \medskip\noindent\textbf{Exactness.}
    
    $\iota$ is injective: $\iota_1=\mathrm{id}_k$ is injective.
    
    $p$ is surjective: $p_2=\mathrm{id}_k$ is surjective.
    
    Exactness at $M_\lambda$: we verify $\ker p = \operatorname{im}\iota$ at each vertex.
    \begin{itemize}
    \item Vertex~$1$:
    \[
    \ker(p_1\colon k\to 0) = k = \operatorname{im}(\iota_1) = \operatorname{im}(\mathrm{id}_k) = k.
    \]
    \item Vertex~$2$:
    \[
    \ker(p_2) = \ker(\mathrm{id}_k) = 0 = \operatorname{im}(\iota_2) = \operatorname{im}(0) = 0.
    \]
    \end{itemize}
    
    Hence
    \[
    0 \longrightarrow S(1) \xrightarrow{\;\iota\;} M_\lambda \xrightarrow{\;p\;} S(2) \longrightarrow 0
    \]
    is a short exact sequence.
    \end{proof}
\end{enumerate}

%=============================================================================
% Problem 7
%=============================================================================
\item Let $f \colon M \to N$ be a morphism of representations. Prove that the definition of $\operatorname{coker} f$ is equivalent to the one defined in classroom through universal property (otherwise, see Remark~1.6 in Schiffler's book).

\begin{proof}
    Let
    \[
    f=(f_i):M=(M_i,\varphi_\alpha)\to N=(N_i,\psi_\alpha)
    \]
    be a morphism of representations. For each vertex $i$, define
    \[
    H_i:=N_i/f_i(M_i).
    \]
    For each arrow $\alpha:i\to j$, define
    \[
    \chi_\alpha\bigl(n_i+f_i(M_i)\bigr):=\psi_\alpha(n_i)+f_j(M_j).
    \]
    
    Hence
    \[
    H=(H_i,\chi_\alpha)
    \]
    is a representation.
    
    Now let
    \[
    g=(g_i):N\to H
    \]
    be the natural quotient morphism, where
    \[
    g_i:N_i\to N_i/f_i(M_i),\qquad g_i(n_i)=n_i+f_i(M_i).
    \]
    Then for every vertex $i$ and every $m_i\in M_i$,
    \[
    (g_i\circ f_i)(m_i)=f_i(m_i)+f_i(M_i)=0.
    \]
    Thus
    \[
    g\circ f=0.
    \]
    
    We now prove the universal property. Let
    \[
    h=(h_i):N\to Y=(Y_i,\eta_\alpha)
    \]
    be a morphism such that
    \[
    h\circ f=0.
    \]
    Then for each vertex $i$,
    \[
    h_i\circ f_i=0.
    \]
    Hence $h_i$ vanishes on $f_i(M_i)$, so it factors uniquely through the quotient
    $N_i/f_i(M_i)$. Therefore there exists a unique linear map
    \[
    u_i:H_i=N_i/f_i(M_i)\to Y_i
    \]
    such that
    \[
    u_i(n_i+f_i(M_i)):=h_i(n_i).
    \]
    This is well-defined because if
    \[
    n_i+f_i(M_i)=n_i'+f_i(M_i),
    \]
    then $n_i-n_i'=f_i(m_i)$ for some $m_i\in M_i$, and hence
    \[
    h_i(n_i)-h_i(n_i')=h_i(f_i(m_i))=0.
    \]
    
    By construction,
    \[
    u_i\circ g_i=h_i
    \]
    for every vertex $i$.
    
    It remains to check that $u=(u_i):H\to Y$ is a morphism of representations.
    Let $\alpha:i\to j$ be an arrow and let $n_i+f_i(M_i)\in H_i$. Then
    \begin{align*}
    \eta_\alpha\bigl(u_i(n_i+f_i(M_i))\bigr)
    &=\eta_\alpha(h_i(n_i))
    =h_j(\psi_\alpha(n_i))
    =u_j\bigl(\psi_\alpha(n_i)+f_j(M_j)\bigr)\\
    &=u_j\bigl(\chi_\alpha(n_i+f_i(M_i))\bigr).
    \end{align*}
    Therefore
    \[
    \eta_\alpha\circ u_i=u_j\circ \chi_\alpha,
    \]
    so $u$ is a morphism of representations.
    
    Finally, uniqueness is clear: if
    \[
    u'=(u_i'):H\to Y
    \]
    also satisfies
    \[
    u'\circ g=h,
    \]
    then for each $i$,
    \[
    u_i'(n_i+f_i(M_i))
    =u_i'(g_i(n_i))
    =h_i(n_i)
    =u_i(n_i+f_i(M_i)).
    \]
    Thus $u_i'=u_i$ for all $i$, so $u'=u$.
    
    Therefore the quotient representation
    \[
    \bigl(N_i/f_i(M_i),\chi_\alpha\bigr)
    \]
    with the natural projection $g:N\to H$ satisfies the universal property of the cokernel.
  
        On the other hand, suppose
        \[
        g=(g_i):N=(N_i,\psi_\alpha)\to H=(H_i,\chi_\alpha)
        \]
        is a cokernel of
        \[
        f=(f_i):M=(M_i,\varphi_\alpha)\to N=(N_i,\psi_\alpha)
        \]
        in the sense of the universal property.
        
        Let
        \[
        C:=\bigl(N_i/f_i(M_i),\ \overline{\psi}_\alpha\bigr)
        \]
        be the pointwise cokernel representation, where
        \[
        \overline{\psi}_\alpha(n_i+f_i(M_i))
        :=\psi_\alpha(n_i)+f_j(M_j).
        \]
        Let
        \[
        \pi=(\pi_i):N\to C
        \]
        be the natural quotient morphism. As shown above,
        \[
        \pi\circ f=0.
        \]
        
        Since $g:N\to H$ is a cokernel of $f$, applying its universal property to the morphism
        \[
        \pi:N\to C
        \]
        with $\pi\circ f=0$, there exists a unique morphism
        \[
        a:H\to C
        \]
        such that
        \[
        a\circ g=\pi.
        \]
        
        On the other hand, since $g\circ f=0$, for each vertex $i$ we have
        \[
        g_i\circ f_i=0.
        \]
        Thus $g_i$ vanishes on $f_i(M_i)$, so it factors uniquely through the quotient
        \[
        N_i/f_i(M_i).
        \]
        Hence for each $i$ there exists a unique linear map
        \[
        b_i:C_i=N_i/f_i(M_i)\to H_i
        \]
        given by
        \[
        b_i(n_i+f_i(M_i)):=g_i(n_i).
        \]
        These maps assemble into a morphism
        \[
        b:C\to H,
        \]
        because for any arrow $\alpha:i\to j$,
        \begin{align*}
        \chi_\alpha\bigl(b_i(n_i+f_i(M_i))\bigr)
        &=\chi_\alpha(g_i(n_i))\\
        &=g_j(\psi_\alpha(n_i))\\
        &=b_j\bigl(\psi_\alpha(n_i)+f_j(M_j)\bigr)\\
        &=b_j\bigl(\overline{\psi}_\alpha(n_i+f_i(M_i))\bigr).
        \end{align*}
        So
        \[
        b\circ \pi=g.
        \]
        
        We now show that $a$ and $b$ are inverse isomorphisms.
        
        First,
        \[
        (a\circ b)\circ \pi
        =a\circ (b\circ \pi)
        =a\circ g
        =\pi.
        \]
        Also,
        \[
        \operatorname{id}_C\circ \pi=\pi.
        \]
        Since $\pi:N\to C$ is a cokernel of $f$, the uniqueness part of its universal property implies
        \[
        a\circ b=\operatorname{id}_C.
        \]
        
        Similarly,
        \[
        (b\circ a)\circ g
        =b\circ (a\circ g)
        =b\circ \pi
        =g,
        \]
        and also
        \[
        \operatorname{id}_H\circ g=g.
        \]
        Since $g:N\to H$ is a cokernel of $f$, the uniqueness part of its universal property implies
        \[
        b\circ a=\operatorname{id}_H.
        \]
        
        Therefore $a$ and $b$ are inverse isomorphisms, and hence
        \[
        H\cong C=\bigl(N_i/f_i(M_i),\ \overline{\psi}_\alpha\bigr).
        \]
        
        In particular, for each vertex $i$ we have
        \[
        H_i\cong N_i/f_i(M_i),
        \]
        and under these identifications, for every arrow $\alpha:i\to j$,
        \[
        \chi_\alpha
        =
        a_j^{-1}\circ \overline{\psi}_\alpha\circ a_i.
        \]
        Thus any cokernel given by the universal property is isomorphic to the pointwise quotient representation
        \[
        \bigl(N_i/f_i(M_i),\ \psi_\alpha \text{ induced on the quotients}\bigr).
        \]
\end{proof}

%=============================================================================
% Problem 8
%=============================================================================
\item Let $M$, $M'$ and $N$ be representations of $Q$ and let $f \colon M \to N$, $g \colon M' \to N$ be morphisms. Define the \emph{fiber product} (or \emph{pull back}) of $f$ and $g$ as
\[
    X = \{(a, b) \mid a \in M,\; b \in M',\; \text{such that } f(a) = g(b)\},
\]
and define the projections $\pi_1(a, b) = a$ and $\pi_2(a, b) = b$.

\begin{enumerate}[label=(\alph*)]
    \item Show that $X$ is a subrepresentation of $M \oplus M'$ and that there is a commutative diagram:
    \[
    \begin{tikzcd}
        X \arrow[r, "\pi_2"] \arrow[d, "\pi_1"'] & M' \arrow[d, "g"] \\
        M \arrow[r, "f"'] & N
    \end{tikzcd}
    \]
    \begin{proof}
    \textbf{$X$ is a subrepresentation of $M\oplus M'$.}
    
    At each vertex $i$, define
    \[
    X_i = \{(a,b)\in M_i\oplus M'_i : f_i(a) = g_i(b)\}.
    \]
    This is the kernel of the linear map
    \[
    M_i\oplus M'_i \longrightarrow N_i,\qquad (a,b)\longmapsto f_i(a)-g_i(b),
    \]
    so $X_i$ is a linear subspace of $M_i\oplus M'_i$.

    We now show that for any arrow $\alpha\colon i\to j$, the map $(M\oplus M')(\alpha)$ sends $X_i$ into $X_j$. Let $(a,b)\in X_i$, so
    \[
    f_i(a)=g_i(b).
    \]
    We need to show $(M(\alpha)(a),\; M'(\alpha)(b)) \in X_j$, i.e., $f_j(M(\alpha)(a)) = g_j(M'(\alpha)(b))$. Using the commutativity of $f$ and $g$:
    \[
    f_j\bigl(M(\alpha)(a)\bigr) = N(\alpha)\bigl(f_i(a)\bigr) = N(\alpha)\bigl(g_i(b)\bigr) = g_j\bigl(M'(\alpha)(b)\bigr).
    \]
    Hence $X$ is stable under the maps of $M\oplus M'$ and is therefore a subrepresentation.

    \medskip\noindent\textbf{Commutativity of the diagram.}
    
    For any $(a,b)\in X_i$, we have
    \[
    (f\circ\pi_1)_i(a,b) = f_i\bigl(\pi_1(a,b)\bigr) = f_i(a) = g_i(b) = g_i\bigl(\pi_2(a,b)\bigr) = (g\circ\pi_2)_i(a,b).
    \]
    Hence $f\circ\pi_1 = g\circ\pi_2$, and the diagram
    \[
    \begin{tikzcd}
        X \arrow[r, "\pi_2"] \arrow[d, "\pi_1"'] & M' \arrow[d, "g"] \\
        M \arrow[r, "f"'] & N
    \end{tikzcd}
    \]
    commutes.
    \end{proof}
    \item Show that if $f$ is injective then $\pi_2$ is injective.
    \begin{proof}
    Suppose $f$ is injective, i.e., $f_i$ is injective for each vertex $i$.
    
    Let $(a,b)\in X_i$ with
    \[
    \pi_2(a,b) = b = 0.
    \]
    Since $(a,b)\in X_i$, we have $f_i(a)=g_i(b)$. Substituting $b=0$:
    \[
    f_i(a) = g_i(0) = 0.
    \]
    Since $f_i$ is injective, it follows that
    \[
    a = 0.
    \]
    Hence $(a,b)=(0,0)$, so $\ker(\pi_2)_i = 0$ at every vertex $i$. Therefore $\pi_2$ is injective.
    \end{proof}
    \item Show that if $f$ is surjective then $\pi_2$ is surjective.
    \begin{proof}
    Suppose $f$ is surjective, i.e., $f_i$ is surjective for each vertex $i$.
    
    Let $b\in M'_i$. We need to find $a\in M_i$ such that $(a,b)\in X_i$, i.e., such that
    \[
    f_i(a) = g_i(b).
    \]
    The element $g_i(b)\in N_i$. Since $f_i$ is surjective, we have
    \[
    N_i = \operatorname{im}(f_i),
    \]
    so there exists $a\in M_i$ with $f_i(a) = g_i(b)$. Then $(a,b)\in X_i$ by definition, and
    \[
    \pi_2(a,b) = b.
    \]
    Since $b\in M'_i$ was arbitrary, $\pi_2$ is surjective at each vertex.
    \end{proof}
    \item Suppose $0 \to L \xrightarrow{h} M \xrightarrow{f} N \to 0$ is a short exact sequence, define $h' \colon L \to X$ by $h'(n) = (h(n), 0)$. Show that the following diagram is commutative with exact rows:
    \[
    \begin{tikzcd}
        0 \arrow[r] & L \arrow[r, "h'"] \arrow[d, "1_L"'] & X \arrow[r, "\pi_2"] \arrow[d, "\pi_1"'] & M' \arrow[r] \arrow[d, "g"] & 0 \\
        0 \arrow[r] & L \arrow[r, "h"'] & M \arrow[r, "f"'] & N \arrow[r] & 0
    \end{tikzcd}
    \]
    \begin{proof}
    \textbf{$h'$ is well-defined.}
    
    For $n\in L_i$, we must check that $(h_i(n),\,0)\in X_i$, i.e., that
    \[
    f_i(h_i(n)) = g_i(0).
    \]
    By exactness of the bottom row at $M$, we have $f\circ h = 0$, so
    \[
    f_i(h_i(n)) = 0 = g_i(0).
    \]
    Hence $h'_i(n) = (h_i(n),\,0) \in X_i$.

    \medskip\noindent\textbf{$h'$ is a morphism of representations.}
    
    For each arrow $\alpha\colon i\to j$ and $n\in L_i$:
    \begin{align*}
    X(\alpha)\bigl(h'_i(n)\bigr) &= X(\alpha)\bigl(h_i(n),\,0\bigr) \\
    &= \bigl(M(\alpha)(h_i(n)),\; M'(\alpha)(0)\bigr) \\
    &= \bigl(h_j(L(\alpha)(n)),\; 0\bigr) \\
    &= h'_j\bigl(L(\alpha)(n)\bigr),
    \end{align*}
    where the third equality uses the commutativity of $h$, i.e., $M(\alpha)\circ h_i = h_j\circ L(\alpha)$.

    \medskip\noindent\textbf{Commutativity of the diagram.}
    
    Left square: for each $n\in L_i$,
    \[
    (\pi_1\circ h')_i(n) = \pi_1(h_i(n),\,0) = h_i(n) = (h\circ 1_L)_i(n).
    \]
    So $\pi_1\circ h' = h\circ 1_L$.
    
    Right square: $f\circ\pi_1 = g\circ\pi_2$ was established in part~(a).

    \medskip\noindent\textbf{Exactness of the top row.}

    \emph{$h'$ is injective:} suppose
    \[
    h'_i(n) = (h_i(n),\,0) = (0,\,0).
    \]
    Then $h_i(n) = 0$. Since $h$ is injective (from the exactness of the bottom row),
    \[
    n = 0.
    \]
    Hence $h'$ is injective.

    \emph{$\pi_2$ is surjective:} since $f$ is surjective (from the bottom row), this follows from part~(c).

    \emph{$\operatorname{im}(h') = \ker(\pi_2)$:}

    $(\supseteq)$ For any $n\in L_i$,
    \[
    \pi_2(h'_i(n)) = \pi_2(h_i(n),\,0) = 0,
    \]
    so $\operatorname{im}(h'_i)\subseteq\ker(\pi_2)_i$.
    
    $(\subseteq)$ Let $(a,b)\in X_i$ with $\pi_2(a,b)=b=0$. Then
    \[
    f_i(a) = g_i(0) = 0,
    \]
    so $a\in\ker f_i$. By exactness of the bottom row at $M$,
    \[
    \ker f_i = \operatorname{im}\, h_i.
    \]
    Hence there exists $n\in L_i$ with $a=h_i(n)$, and then
    \[
    (a,\,0) = (h_i(n),\,0) = h'_i(n) \in \operatorname{im}(h'_i).
    \]
    
    Therefore $\operatorname{im}(h') = \ker(\pi_2)$, and the top row is exact.
    \end{proof}
\end{enumerate}

%=============================================================================
% Problem 9
%=============================================================================
\item Let $f \colon L \to M$ be a morphism. Then a \emph{cokernel} of $f$ is a morphism $g \colon M \to N$ such that $g f = 0$, and given any morphism $v \colon M \to X$ such that $v f = 0$, there is a unique morphism $u \colon N \to X$ such that $u g = v$:
\[
\begin{tikzcd}
L \arrow[r, "f"] & M \arrow[r, "g"] \arrow[d, "v"'] & N \arrow[dl, dashed, "\exists\, u"] \\
& X
\end{tikzcd}
\]
Prove that the two definitions of cokernel agree.

\begin{proof}
    Let
    \[
    f=(f_i)\colon L=(L_i,\varphi_\alpha)\to M=(M_i,\psi_\alpha)
    \]
    be a morphism of representations.
    
    We first show that the pointwise cokernel satisfies the universal property. For each vertex $i\in Q_0$, set
    \[
    H_i:=M_i/f_i(L_i),
    \]
    and let
    \[
    \pi_i\colon M_i\to H_i
    \]
    be the canonical quotient map. For an arrow $\alpha\colon i\to j$, define
    \[
    \bar{\psi}_\alpha\colon H_i\to H_j
    \]
    by
    \[
    \bar{\psi}_\alpha\bigl(m_i+f_i(L_i)\bigr):=\psi_\alpha(m_i)+f_j(L_j).
    \]
    This is well defined because $f$ is a morphism, so
    \[
    \psi_\alpha\bigl(f_i(L_i)\bigr)=f_j\bigl(\varphi_\alpha(L_i)\bigr)\subseteq f_j(L_j).
    \]
    Hence
    \[
    H=(H_i,\bar{\psi}_\alpha)
    \]
    is a representation, and $\pi=(\pi_i)\colon M\to H$ is a morphism.
    
    Moreover,
    \[
    (\pi\circ f)_i=\pi_i\circ f_i=0
    \]
    for every vertex $i$, since $f_i(L_i)\subseteq \ker \pi_i$. Thus
    \[
    \pi\circ f=0.
    \]
    
    Now let $v=(v_i)\colon M\to X=(X_i,\eta_\alpha)$ be a morphism such that
    \[
    v\circ f=0.
    \]
    Then for every vertex $i$ we have
    \[
    v_i\circ f_i=0,
    \]
    so $v_i$ vanishes on $f_i(L_i)$. By the universal property of the quotient vector space $H_i=M_i/f_i(L_i)$, there exists a unique linear map
    \[
    u_i\colon H_i\to X_i
    \]
    such that
    \[
    u_i\circ \pi_i=v_i.
    \]
    We claim that the family $(u_i)$ defines a morphism of representations
    \[
    u\colon H\to X.
    \]
    Indeed, let $\alpha\colon i\to j$ be an arrow. Then
    \[
    \eta_\alpha\circ u_i\circ \pi_i
    =
    \eta_\alpha\circ v_i
    =
    v_j\circ \psi_\alpha
    =
    u_j\circ \pi_j\circ \psi_\alpha
    =
    u_j\circ \bar{\psi}_\alpha\circ \pi_i.
    \]
    Since $\pi_i$ is surjective, it follows that
    \[
    \eta_\alpha\circ u_i=u_j\circ \bar{\psi}_\alpha.
    \]
    Therefore $u\colon H\to X$ is a morphism, and by construction
    \[
    u\circ \pi=v.
    \]
    
    To prove uniqueness, suppose $u'\colon H\to X$ is another morphism such that
    \[
    u'\circ \pi=v.
    \]
    Then for every vertex $i$,
    \[
    u'_i\circ \pi_i=v_i=u_i\circ \pi_i.
    \]
    Since $\pi_i$ is surjective, we get $u'_i=u_i$. Hence $u'=u$. Thus $\pi\colon M\to H$ satisfies the universal property of the cokernel of $f$.
    
    Next, suppose that
    \[
    g\colon M\to N
    \]
    is any cokernel of $f$ in the sense of the universal property. Since $\pi\circ f=0$, the universal property of $g$ yields a unique morphism
    \[
    a\colon N\to H
    \]
    such that
    \[
    a\circ g=\pi.
    \]
    Similarly, since $g\circ f=0$, the universal property of $\pi$ yields a unique morphism
    \[
    b\colon H\to N
    \]
    such that
    \[
    b\circ \pi=g.
    \]
    
    We now show that $a$ and $b$ are inverse isomorphisms. First,
    \[
    (a\circ b)\circ \pi
    =
    a\circ (b\circ \pi)
    =
    a\circ g
    =
    \pi
    =
    \operatorname{id}_H\circ \pi.
    \]
    Both $a\circ b$ and $\operatorname{id}_H$ are morphisms $H\to H$ whose composition with $\pi$ is equal to $\pi$. By the universal property of $\pi$, such a morphism is unique. Therefore
    \[
    a\circ b=\operatorname{id}_H.
    \]
    
    Similarly,
    \[
    (b\circ a)\circ g
    =
    b\circ (a\circ g)
    =
    b\circ \pi
    =
    g
    =
    \operatorname{id}_N\circ g.
    \]
    Both $b\circ a$ and $\operatorname{id}_N$ are morphisms $N\to N$ whose composition with $g$ is equal to $g$. By the universal property of $g$, such a morphism is unique. Hence
    \[
    b\circ a=\operatorname{id}_N.
    \]
    
    Therefore $a$ and $b$ are inverse isomorphisms, so $N\cong H$. This proves that any cokernel in the universal-property sense is isomorphic to the pointwise cokernel. Hence the two definitions agree.
    \end{proof}

%=============================================================================
% Problem 10
%=============================================================================
\item Let $Q$ be a quiver and let
\[
L \xrightarrow{\;f\;} M \xrightarrow{\;g\;} N \longrightarrow 0
\]
be a sequence in $\operatorname{Rep} Q$. Show that this sequence is exact if and only if for every representation $X \in \operatorname{Rep} Q$, the sequence
\[
0 \longrightarrow \operatorname{Hom}(N, X) \xrightarrow{\;g^\ast\;} \operatorname{Hom}(M, X) \xrightarrow{\;f^\ast\;} \operatorname{Hom}(L, X)
\]
is exact, where $g^\ast(h) = h \circ g$ and $f^\ast(k) = k \circ f$.

\begin{proof}
    $(\Rightarrow)$ Assume that
    \[
    L \xrightarrow{f} M \xrightarrow{g} N \longrightarrow 0
    \]
    is exact. Let $X \in \operatorname{Rep}Q$ be arbitrary. We show that
    \[
    0 \longrightarrow \operatorname{Hom}(N,X)
    \xrightarrow{\,g^\ast\,}
    \operatorname{Hom}(M,X)
    \xrightarrow{\,f^\ast\,}
    \operatorname{Hom}(L,X)
    \]
    is exact.
    
    First, $g^\ast$ is injective. Indeed, if $h,h' \in \operatorname{Hom}(N,X)$ satisfy
    \[
    g^\ast(h)=g^\ast(h'),
    \]
    then
    \[
    h\circ g=h'\circ g.
    \]
    Since $g$ is surjective, it follows that $h=h'$.
    
    Next, for any $h \in \operatorname{Hom}(N,X)$,
    \[
    f^\ast(g^\ast(h))=(h\circ g)\circ f=h\circ(g\circ f)=0,
    \]
    because $g\circ f=0$. Hence
    \[
    \operatorname{im} g^\ast \subseteq \ker f^\ast.
    \]
    
    Conversely, let $k \in \operatorname{Hom}(M,X)$ satisfy
    \[
    f^\ast(k)=k\circ f=0.
    \]
    Since the sequence
    \[
    L \xrightarrow{f} M \xrightarrow{g} N \longrightarrow 0
    \]
    is exact, the morphism $g$ is a cokernel of $f$. Therefore, by the universal
    property of the cokernel, there exists a unique morphism
    \[
    u:N\to X
    \]
    such that
    \[
    u\circ g=k.
    \]
    Thus $k=g^\ast(u)$, and so
    \[
    \ker f^\ast \subseteq \operatorname{im} g^\ast.
    \]
    Therefore
    \[
    \ker f^\ast=\operatorname{im} g^\ast,
    \]
    and the induced sequence is exact.
    
    $(\Leftarrow)$ Assume that for every $X \in \operatorname{Rep}Q$, the sequence
    \[
    0 \longrightarrow \operatorname{Hom}(N,X)
    \xrightarrow{\,g^\ast\,}
    \operatorname{Hom}(M,X)
    \xrightarrow{\,f^\ast\,}
    \operatorname{Hom}(L,X)
    \]
    is exact. We show that
    \[
    L \xrightarrow{f} M \xrightarrow{g} N \longrightarrow 0
    \]
    is exact.
    
    First, we prove that $g\circ f=0$. Take $X=N$. Since
    \[
    f^\ast\circ g^\ast=0,
    \]
    evaluating at $\operatorname{id}_N \in \operatorname{Hom}(N,N)$ gives
    \[
    0=f^\ast(g^\ast(\operatorname{id}_N))
    =\operatorname{id}_N\circ g\circ f
    =g\circ f.
    \]
    
    Next, we show that $g$ is surjective. Let
    \[
    \pi:N\to \operatorname{coker} g
    \]
    be the cokernel morphism of $g$. Since $\pi\circ g=0$, we have
    \[
    g^\ast(\pi)=\pi\circ g=0.
    \]
    By exactness, $g^\ast$ is injective, hence $\pi=0$. Therefore
    \[
    \operatorname{coker} g=0,
    \]
    so $g$ is surjective.
    
    Finally, we show that $\ker g=\operatorname{im} f$. Since $g\circ f=0$, we have
    \[
    \operatorname{im} f \subseteq \ker g.
    \]
    For the reverse inclusion, let
    \[
    \gamma:M\to \operatorname{coker} f
    \]
    be the cokernel morphism of $f$. Since $\gamma\circ f=0$, we have
    \[
    f^\ast(\gamma)=0.
    \]
    By exactness of
    \[
    0 \longrightarrow \operatorname{Hom}(N,\operatorname{coker}f)
    \xrightarrow{\,g^\ast\,}
    \operatorname{Hom}(M,\operatorname{coker}f)
    \xrightarrow{\,f^\ast\,}
    \operatorname{Hom}(L,\operatorname{coker}f),
    \]
    it follows that $\gamma \in \operatorname{im} g^\ast$. Hence there exists a morphism
    \[
    u:N\to \operatorname{coker} f
    \]
    such that
    \[
    \gamma=u\circ g.
    \]
    Therefore
    \[
    \ker g \subseteq \ker \gamma.
    \]
    Since $\gamma$ is the cokernel morphism of $f$, we have
    \[
    \ker \gamma=\operatorname{im} f.
    \]
    Thus
    \[
    \ker g \subseteq \operatorname{im} f.
    \]
    Combining this with the reverse inclusion, we get
    \[
    \ker g=\operatorname{im} f.
    \]
    Hence
    \[
    L \xrightarrow{f} M \xrightarrow{g} N \longrightarrow 0
    \]
    is exact.
    \end{proof}

%=============================================================================
% Problem 11
%=============================================================================
\item Let $g \colon M \to N$ be a surjective morphism between representations of $Q$, and let $P(i)$ be the projective representation at vertex~$i$. Show that the map
\[
g_\ast \colon \operatorname{Hom}(P(i), M) \to \operatorname{Hom}(P(i), N),\qquad \varphi \longmapsto g \circ \varphi
\]
is surjective.

\begin{proof}
    Let
    \[
    g_\ast \colon \operatorname{Hom}(P(i),M)\longrightarrow \operatorname{Hom}(P(i),N),
    \qquad
    \varphi \longmapsto g\circ \varphi .
    \]
    We must show that \(g_\ast\) is surjective.
    
    Recall that for every representation \(X\in \operatorname{Rep}Q\), there is a natural isomorphism
    \[
    \Phi_X \colon \operatorname{Hom}(P(i),X)\xrightarrow{\;\sim\;} X_i,
    \qquad
    f\longmapsto f_i(e_i),
    \]
    where \(e_i\) denotes the trivial path at the vertex \(i\).
    
    We claim that the following diagram commutes:
    \[
    \begin{tikzcd}
    \operatorname{Hom}(P(i),M) \arrow[r, "\Phi_M", "\sim"'] \arrow[d, "g_\ast"'] 
    & M_i \arrow[d, "g_i"] \\
    \operatorname{Hom}(P(i),N) \arrow[r, "\Phi_N", "\sim"'] 
    & N_i
    \end{tikzcd}
    \]
    Indeed, for any \(f\in \operatorname{Hom}(P(i),M)\), we have
    \[
    \Phi_N\bigl(g_\ast(f)\bigr)
    =
    \Phi_N(g\circ f)
    =
    (g\circ f)_i(e_i)
    =
    g_i\bigl(f_i(e_i)\bigr)
    =
    g_i\bigl(\Phi_M(f)\bigr).
    \]
    Thus
    \[
    \Phi_N\circ g_\ast = g_i\circ \Phi_M.
    \]
    
    Since \(g\colon M\to N\) is surjective, each component map
    \[
    g_j\colon M_j\to N_j
    \]
    is surjective. In particular, \(g_i\colon M_i\to N_i\) is surjective. Because \(\Phi_M\) and \(\Phi_N\) are isomorphisms, the commutative diagram above shows that \(g_\ast\) corresponds to the surjective map \(g_i\). Hence \(g_\ast\) is surjective.
    
    Therefore every morphism \(P(i)\to N\) lifts to a morphism \(P(i)\to M\), as required.
    \end{proof}

%=============================================================================
% Problem 12
%=============================================================================
\item Let $g \colon L \to M$ be an injective morphism between representations of $Q$, and let $I(i)$ be the injective representation at vertex~$i$. Show that the map
\[
g^\ast \colon \operatorname{Hom}(M, I(i)) \to \operatorname{Hom}(L, I(i)),\qquad \psi \longmapsto \psi \circ g
\]
is surjective.

\begin{proof}
    For every representation \(X\), there is a natural isomorphism
    \[
    \Psi_X \colon \operatorname{Hom}(X,I(i)) \xrightarrow{\;\sim\;} X_i^{\ast},
    \]
    where \(X_i^{\ast}=\operatorname{Hom}_k(X_i,k)\). Concretely, after identifying
    \[
    I(i)_i \cong k
    \]
    via the trivial path \(e_i\), a morphism \(f\colon X\to I(i)\) is sent to the linear functional
    \[
    \Psi_X(f)=f_i \in \operatorname{Hom}_k(X_i,k)=X_i^{\ast}.
    \]
    
    Under these identifications, the map
    \[
    g^\ast \colon \operatorname{Hom}(M,I(i)) \to \operatorname{Hom}(L,I(i)), 
    \qquad \psi \mapsto \psi\circ g,
    \]
    corresponds to the dual map
    \[
    g_i^{\ast} \colon M_i^{\ast} \to L_i^{\ast}, 
    \qquad \lambda \mapsto \lambda\circ g_i.
    \]
    Indeed, for every \(\psi \in \operatorname{Hom}(M,I(i))\), we have
    \[
    \Psi_L\bigl(g^\ast(\psi)\bigr)
     = \Psi_L(\psi\circ g)
     = (\psi\circ g)_i
     = \psi_i\circ g_i
     = g_i^{\ast}(\psi_i)
     = g_i^{\ast}\bigl(\Psi_M(\psi)\bigr).
    \]
    Hence the following diagram commutes:
    \[
    \begin{tikzcd}
    \operatorname{Hom}(M,I(i)) \arrow[r, "\Psi_M", "\sim"'] \arrow[d, "g^\ast"'] 
    & M_i^\ast \arrow[d, "g_i^\ast"] \\
    \operatorname{Hom}(L,I(i)) \arrow[r, "\Psi_L", "\sim"'] 
    & L_i^\ast
    \end{tikzcd}
    \]
    
    Now \(g\colon L\to M\) is injective, so each component
    \[
    g_j\colon L_j\to M_j
    \]
    is injective; in particular, \(g_i\) is injective. Since we are working with finite-dimensional representations, \(L_i\) and \(M_i\) are finite-dimensional vector spaces. Therefore the dual map
    \[
    g_i^\ast \colon M_i^\ast \to L_i^\ast
    \]
    is surjective.
    
    For completeness, let \(\lambda \in L_i^\ast\). Since \(g_i\) is injective, it induces an isomorphism
    \[
    L_i \xrightarrow{\;\sim\;} g_i(L_i)\subseteq M_i.
    \]
    Thus \(\lambda\) defines a linear form on the subspace \(g_i(L_i)\) by
    \[
    g_i(x)\longmapsto \lambda(x).
    \]
    By linear algebra, this linear form extends to some \(\mu \in M_i^\ast\). Then
    \[
    g_i^\ast(\mu)=\mu\circ g_i=\lambda.
    \]
    Hence \(g_i^\ast\) is surjective.
    
    Since \(\Psi_M\) and \(\Psi_L\) are isomorphisms and the diagram above commutes, it follows that
    \[
    g^\ast \colon \operatorname{Hom}(M,I(i)) \to \operatorname{Hom}(L,I(i))
    \]
    is surjective.
    \end{proof}

%=============================================================================
% Problem 13
%=============================================================================
\item Let $Q$ be the quiver
\[
    \begin{tikzcd}
        1 \arrow[rr] \arrow[rd] &              & 3 \arrow[r] & 4 \\
                                & 2 \arrow[ru] &             &  
        \end{tikzcd}
\]
Prove that
\[
P(1)\simeq
\begin{tikzcd}[ampersand replacement=\&, column sep=large, row sep=large]
k \arrow[rr, "{\begin{bmatrix}x_1\\x_2\end{bmatrix}}"] \arrow[dr, "z"'] \& \& k^2 \arrow[r, "{\begin{bmatrix}a&b\\ c&d\end{bmatrix}}"] \& k^2 \\
\& k \arrow[ur, "{\begin{bmatrix}y_1\\y_2\end{bmatrix}}"'] \&
\end{tikzcd}
\]
if and only if all four maps have maximal rank, that is,
\[
ad-bc\neq 0,\qquad (x_1,x_2)\neq (0,0),\qquad (y_1,y_2)\neq (0,0),\qquad z\neq 0
\]
and the vectors
\[
\begin{bmatrix}a&b\\ c&d\end{bmatrix}
\begin{bmatrix}x_1\\x_2\end{bmatrix}
\qquad\text{and}\qquad
\begin{bmatrix}a&b\\ c&d\end{bmatrix}
\begin{bmatrix}y_1\\y_2\end{bmatrix}
\]
are linearly independent.

\begin{proof}
    Let the arrows of \(Q\) be
    \[
    \alpha\colon 1\to 3,\qquad
    \beta\colon 1\to 2,\qquad
    \gamma\colon 2\to 3,\qquad
    \delta\colon 3\to 4,
    \]
    and let \(M\) denote the representation
    \[
    \begin{tikzcd}[ampersand replacement=\&, column sep=large, row sep=large]
    k \arrow[rr, "{\begin{bmatrix}x_1\\x_2\end{bmatrix}}"] \arrow[dr, "z"'] \& \& k^2 \arrow[r, "{\begin{bmatrix}a&b\\ c&d\end{bmatrix}}"] \& k^2 \\
    \& k \arrow[ur, "{\begin{bmatrix}y_1\\y_2\end{bmatrix}}"'] \&
    \end{tikzcd}
    \]
    that is,
    \[
    M_\alpha=\begin{bmatrix}x_1\\x_2\end{bmatrix},\qquad
    M_\beta=z,\qquad
    M_\gamma=\begin{bmatrix}y_1\\y_2\end{bmatrix},\qquad
    M_\delta=\begin{bmatrix}a&b\\c&d\end{bmatrix}.
    \]
    
    \medskip
    \noindent\textbf{Step 1: Description of \(P(1)\).}
    Since \(Q\) has no oriented cycles, the projective representation \(P(1)\) is given at each vertex \(j\) by the \(k\)-vector space with basis the set of paths from \(1\) to \(j\). Thus
    \[
    P(1)_1 \cong k,\qquad
    P(1)_2 \cong k,\qquad
    P(1)_3 \cong k^2,\qquad
    P(1)_4 \cong k^2,
    \]
    because:
    \begin{itemize}
        \item there is one path \(1\to 1\), namely the trivial path \(e_1\);
        \item there is one path \(1\to 2\), namely \(\beta\);
        \item there are two paths \(1\to 3\), namely \(\alpha\) and \(\gamma\beta\);
        \item there are two paths \(1\to 4\), namely \(\delta\alpha\) and \(\delta\gamma\beta\).
    \end{itemize}
    
    Choose the ordered bases
    \[
    P(1)_1=(e_1),\qquad
    P(1)_2=(\beta),\qquad
    P(1)_3=(\alpha,\gamma\beta),\qquad
    P(1)_4=(\delta\alpha,\delta\gamma\beta).
    \]
    With respect to these bases, the structure maps of \(P(1)\) are
    \[
    P(1)_\beta=1,\qquad
    P(1)_\alpha=\begin{bmatrix}1\\0\end{bmatrix},\qquad
    P(1)_\gamma=\begin{bmatrix}0\\1\end{bmatrix},\qquad
    P(1)_\delta=I_2.
    \]
    Hence \(P(1)\) has the form
    \[
    P(1)\simeq
    \begin{tikzcd}[ampersand replacement=\&, column sep=large, row sep=large]
    k \arrow[rr, "{\begin{bmatrix}1\\0\end{bmatrix}}"] \arrow[dr, "1"'] \& \& k^2 \arrow[r, "I_2"] \& k^2 \\
    \& k \arrow[ur, "{\begin{bmatrix}0\\1\end{bmatrix}}"'] \&
    \end{tikzcd}
    \]
    In particular, all four maps have maximal rank, and the two vectors
    \[
    P(1)_\delta P(1)_\alpha=\begin{bmatrix}1\\0\end{bmatrix},
    \qquad
    P(1)_\delta P(1)_\gamma=\begin{bmatrix}0\\1\end{bmatrix}
    \]
    are linearly independent.
    
    \medskip
    \noindent\textbf{Step 2: Necessity.}
    Assume that \(M\simeq P(1)\). Let
    \[
    \phi=(\phi_1,\phi_2,\phi_3,\phi_4)\colon P(1)\xrightarrow{\sim} M
    \]
    be an isomorphism of representations.
    
    Since isomorphisms preserve rank, and since the maps of \(P(1)\) all have maximal rank, it follows immediately that the four maps of \(M\) also have maximal rank. Thus
    \[
    ad-bc\neq 0,\qquad
    (x_1,x_2)\neq (0,0),\qquad
    (y_1,y_2)\neq (0,0),\qquad
    z\neq 0.
    \]
    
    It remains to show that
    \[
    M_\delta M_\alpha
    \qquad\text{and}\qquad
    M_\delta M_\gamma
    \]
    are linearly independent in \(k^2\).
    
    Since \(\phi_1,\phi_2\colon k\to k\) are automorphisms, there exist nonzero scalars
    \[
    \phi_1=\lambda_1,\qquad \phi_2=\lambda_2
    \]
    with \(\lambda_1,\lambda_2\in k^\times\). The commutativity of the squares gives
    \[
    \phi_3 P(1)_\alpha=M_\alpha\phi_1,
    \qquad
    \phi_3 P(1)_\gamma=M_\gamma\phi_2,
    \qquad
    \phi_4 P(1)_\delta=M_\delta\phi_3.
    \]
    Therefore
    \[
    M_\delta M_\alpha\,\lambda_1
    =
    M_\delta\phi_3P(1)_\alpha
    =
    \phi_4P(1)_\delta P(1)_\alpha,
    \]
    and similarly
    \[
    M_\delta M_\gamma\,\lambda_2
    =
    \phi_4P(1)_\delta P(1)_\gamma.
    \]
    Since \(\phi_4\) is invertible and \(P(1)_\delta P(1)_\alpha\), \(P(1)_\delta P(1)_\gamma\) are linearly independent, the vectors
    \[
    \phi_4P(1)_\delta P(1)_\alpha
    \qquad\text{and}\qquad
    \phi_4P(1)_\delta P(1)_\gamma
    \]
    are linearly independent. Multiplication by the nonzero scalars \(\lambda_1^{-1}\) and \(\lambda_2^{-1}\) does not affect linear independence, so
    \[
    M_\delta M_\alpha
    \qquad\text{and}\qquad
    M_\delta M_\gamma
    \]
    are linearly independent as required.
    
    \medskip
    \noindent\textbf{Step 3: Sufficiency.}
    Conversely, assume that all four maps of \(M\) have maximal rank, namely
    \[
    ad-bc\neq 0,\qquad
    (x_1,x_2)\neq (0,0),\qquad
    (y_1,y_2)\neq (0,0),\qquad
    z\neq 0,
    \]
    and assume moreover that the two vectors
    \[
    M_\delta M_\alpha=
    \begin{bmatrix}a&b\\c&d\end{bmatrix}\begin{bmatrix}x_1\\x_2\end{bmatrix},
    \qquad
    M_\delta M_\gamma=
    \begin{bmatrix}a&b\\c&d\end{bmatrix}\begin{bmatrix}y_1\\y_2\end{bmatrix}
    \]
    are linearly independent.
    
    Since \(ad-bc\neq 0\), the matrix \(M_\delta\) is invertible. Hence the linear independence of
    \[
    M_\delta\begin{bmatrix}x_1\\x_2\end{bmatrix}
    \qquad\text{and}\qquad
    M_\delta\begin{bmatrix}y_1\\y_2\end{bmatrix}
    \]
    implies that
    \[
    \begin{bmatrix}x_1\\x_2\end{bmatrix}
    \qquad\text{and}\qquad
    \begin{bmatrix}y_1\\y_2\end{bmatrix}
    \]
    are linearly independent. In particular,
    \[
    x_1y_2-x_2y_1\neq 0.
    \]
    
    We now define maps
    \[
    \phi_1=1\colon k\to k,\qquad
    \phi_2=z\colon k\to k,
    \]
    \[
    \phi_3=
    \begin{bmatrix}
    x_1 & zy_1\\
    x_2 & zy_2
    \end{bmatrix}\colon k^2\to k^2,
    \qquad
    \phi_4=M_\delta\phi_3\colon k^2\to k^2.
    \]
    We claim that
    \[
    \phi=(\phi_1,\phi_2,\phi_3,\phi_4)\colon P(1)\to M
    \]
    is an isomorphism of representations.
    
    Indeed, we check the commutativity of the four arrow maps:
    \[
    \phi_2 P(1)_\beta
    =
    z\cdot 1
    =
    z
    =
    M_\beta\cdot 1
    =
    M_\beta\phi_1,
    \]
    \[
    \phi_3 P(1)_\alpha
    =
    \phi_3\begin{bmatrix}1\\0\end{bmatrix}
    =
    \begin{bmatrix}x_1\\x_2\end{bmatrix}
    =
    M_\alpha\phi_1,
    \]
    \[
    \phi_3 P(1)_\gamma
    =
    \phi_3\begin{bmatrix}0\\1\end{bmatrix}
    =
    z\begin{bmatrix}y_1\\y_2\end{bmatrix}
    =
    M_\gamma\phi_2,
    \]
    and
    \[
    \phi_4 P(1)_\delta
    =
    \phi_4 I_2
    =
    \phi_4
    =
    M_\delta\phi_3.
    \]
    Thus \(\phi\) is a morphism of representations.
    
    Finally, each \(\phi_i\) is invertible:
    \[
    \phi_1=1\neq 0,\qquad \phi_2=z\neq 0,
    \]
    \[
    \det(\phi_3)=z(x_1y_2-x_2y_1)\neq 0,
    \]
    and
    \[
    \det(\phi_4)=\det(M_\delta)\det(\phi_3)\neq 0.
    \]
    Hence \(\phi\) is an isomorphism, and therefore
    \[
    M\simeq P(1).
    \]
    
    This proves that
    \[
    P(1)\simeq
    \begin{tikzcd}[ampersand replacement=\&, column sep=large, row sep=large]
    k \arrow[rr, "{\begin{bmatrix}x_1\\x_2\end{bmatrix}}"] \arrow[dr, "z"'] \& \& k^2 \arrow[r, "{\begin{bmatrix}a&b\\ c&d\end{bmatrix}}"] \& k^2 \\
    \& k \arrow[ur, "{\begin{bmatrix}y_1\\y_2\end{bmatrix}}"'] \&
    \end{tikzcd}
    \]
    if and only if all four maps have maximal rank and the two vectors
    \[
    \begin{bmatrix}a&b\\ c&d\end{bmatrix}
    \begin{bmatrix}x_1\\x_2\end{bmatrix},
    \qquad
    \begin{bmatrix}a&b\\ c&d\end{bmatrix}
    \begin{bmatrix}y_1\\y_2\end{bmatrix}
    \]
    are linearly independent.
    \end{proof}

%=============================================================================
% Problem 14
%=============================================================================
\item Compute the indecomposable projective (resp.\ injective) representations $P(i)$ (resp.\ $I(i)$) for the following quivers:
\[
Q_1\colon\quad
\begin{tikzcd}
    1 & 2 \arrow[l, shift left] \arrow[l, shift right] & 3 \arrow[l, shift right] \arrow[l, shift left]
    \end{tikzcd}
\qquad\qquad
Q_2\colon\quad
\begin{tikzcd}
    1 \arrow[r] \arrow[rr, bend left] & 2 \arrow[r] \arrow[d] & 3 \arrow[ld] \\
                                      & 4                     &             
    \end{tikzcd}
\]

\begin{proof}
    We describe the indecomposable projective and injective representations explicitly, up to isomorphism, by choosing natural path bases.
    
    Recall that for a quiver $Q$ and a vertex $i\in Q_0$,
    \[
    P(i)_j \cong k\{\text{paths } i\to j\},
    \qquad
    I(i)_j \cong k\{\text{paths } j\to i\},
    \]
    where $k\{S\}$ denotes the $k$-vector space with basis $S$.
    The linear map attached to an arrow is induced by composition of paths.
    
    \medskip
    
    \noindent
    We use the convention that a product of arrows is read from right to left:
    for instance, $\alpha\gamma$ means ``first $\gamma$, then $\alpha$''.
    
    \bigskip
    \noindent\textbf{(1) The quiver $Q_1$.}
    \[
    Q_1 \colon \quad
    \begin{tikzcd}
    1 & 2 \arrow[l, shift left] \arrow[l, shift right] & 3 \arrow[l, shift left] \arrow[l, shift right]
    \end{tikzcd}
    \]
    Label the arrows by
    \[
    \alpha,\beta \colon 2 \to 1,
    \qquad
    \gamma,\delta \colon 3 \to 2.
    \]
    
    \medskip
    \noindent\textbf{Indecomposable projectives.}
    
    \smallskip
    \noindent\emph{The projective $P(1)$.}
    Since $1$ is a sink, the only path starting at $1$ is the trivial path $e_1$.
    Hence
    \[
    P(1)=S(1)\colon\quad
    \begin{tikzcd}
    k & 0 \arrow[l, shift left] \arrow[l, shift right] & 0 \arrow[l, shift left] \arrow[l, shift right]
    \end{tikzcd}.
    \]
    
    \smallskip
    \noindent\emph{The projective $P(2)$.}
    The paths starting at $2$ are
    \[
    e_2 \text{ (ending at $2$)}, \qquad \alpha,\beta \text{ (ending at $1$)}.
    \]
    Thus
    \[
    P(2)_1 \cong k\{\alpha,\beta\}\cong k^2,\qquad
    P(2)_2 \cong k\{e_2\}\cong k,\qquad
    P(2)_3=0.
    \]
    With respect to the ordered basis $\{\alpha,\beta\}$ at vertex $1$, we obtain
    \[
    P(2)\colon\quad
    \begin{tikzcd}
    k^2 & k \arrow[l, "{\begin{psmallmatrix}1\\0\end{psmallmatrix}}", shift left]
          \arrow[l, "{\begin{psmallmatrix}0\\1\end{psmallmatrix}}"', shift right]
        & 0 \arrow[l, shift left] \arrow[l, shift right]
    \end{tikzcd}.
    \]
    
    \smallskip
    \noindent\emph{The projective $P(3)$.}
    The paths starting at $3$ are
    \[
    e_3 \text{ (ending at $3$)},\qquad
    \gamma,\delta \text{ (ending at $2$)},
    \]
    and
    \[
    \alpha\gamma,\alpha\delta,\beta\gamma,\beta\delta \text{ (ending at $1$)}.
    \]
    Hence
    \[
    P(3)_1 \cong k^4,\qquad P(3)_2 \cong k^2,\qquad P(3)_3 \cong k,
    \]
    where we choose the ordered bases
    \[
    P(3)_1:\ \{\alpha\gamma,\alpha\delta,\beta\gamma,\beta\delta\},
    \qquad
    P(3)_2:\ \{\gamma,\delta\}.
    \]
    Then
    \[
    P(3)\colon\quad
    \begin{tikzcd}[ampersand replacement=\&]
    k^4
    \&
    k^2 \arrow[l, "{\begin{psmallmatrix}
    1&0\\
    0&1\\
    0&0\\
    0&0
    \end{psmallmatrix}}", shift left]
    \arrow[l, "{\begin{psmallmatrix}
    0&0\\
    0&0\\
    1&0\\
    0&1
    \end{psmallmatrix}}"', shift right]
    \&
    k \arrow[l, "{\begin{psmallmatrix}1\\0\end{psmallmatrix}}", shift left]
      \arrow[l, "{\begin{psmallmatrix}0\\1\end{psmallmatrix}}"', shift right]
    \end{tikzcd}.
    \]
    
    \medskip
    \noindent\textbf{Indecomposable injectives.}
    
    \smallskip
    \noindent\emph{The injective $I(3)$.}
    Since $3$ is a source, the only path ending at $3$ is the trivial path $e_3$.
    Therefore
    \[
    I(3)=S(3)\colon\quad
    \begin{tikzcd}
    0 & 0 \arrow[l, shift left] \arrow[l, shift right] & k \arrow[l, shift left] \arrow[l, shift right]
    \end{tikzcd}.
    \]
    
    \smallskip
    \noindent\emph{The injective $I(2)$.}
    The paths ending at $2$ are
    \[
    e_2 \text{ (starting at $2$)}, \qquad \gamma,\delta \text{ (starting at $3$)}.
    \]
    Thus
    \[
    I(2)_1=0,\qquad I(2)_2\cong k,\qquad I(2)_3\cong k^2.
    \]
    With respect to the ordered basis $\{\gamma,\delta\}$ at vertex $3$, we get
    \[
    I(2)\colon\quad
    \begin{tikzcd}
    0 & k \arrow[l, shift left] \arrow[l, shift right]
      & k^2 \arrow[l, "{(1,0)}", shift left]
            \arrow[l, "{(0,1)}"', shift right]
    \end{tikzcd}.
    \]
    
    \smallskip
    \noindent\emph{The injective $I(1)$.}
    The paths ending at $1$ are
    \[
    e_1 \text{ (starting at $1$)},\qquad
    \alpha,\beta \text{ (starting at $2$)},
    \]
    and
    \[
    \alpha\gamma,\alpha\delta,\beta\gamma,\beta\delta \text{ (starting at $3$)}.
    \]
    Hence
    \[
    I(1)_1\cong k,\qquad I(1)_2\cong k^2,\qquad I(1)_3\cong k^4.
    \]
    Choose the ordered bases
    \[
    I(1)_2:\ \{\alpha,\beta\},
    \qquad
    I(1)_3:\ \{\alpha\gamma,\alpha\delta,\beta\gamma,\beta\delta\}.
    \]
    Then
    \[
    I(1)\colon\quad
    \begin{tikzcd}[ampersand replacement=\&]
    k
    \&
    k^2 \arrow[l, "{(1,0)}", shift left]
        \arrow[l, "{(0,1)}"', shift right]
    \&
    k^4 \arrow[l, "{\begin{psmallmatrix}
    1&0&0&0\\
    0&0&1&0
    \end{psmallmatrix}}", shift left]
        \arrow[l, "{\begin{psmallmatrix}
    0&1&0&0\\
    0&0&0&1
    \end{psmallmatrix}}"', shift right]
    \end{tikzcd}.
    \]
    
    \bigskip
    \noindent\textbf{(2) The quiver $Q_2$.}
    \[
    Q_2 \colon \quad
    \begin{tikzcd}
    1 \arrow[r] \arrow[rr, bend left] & 2 \arrow[r] \arrow[d] & 3 \arrow[ld] \\
                                      & 4 &
    \end{tikzcd}
    \]
    Label the arrows by
    \[
    \alpha\colon 1\to 2,\qquad
    \beta\colon 1\to 3,\qquad
    \gamma\colon 2\to 3,\qquad
    \delta\colon 2\to 4,\qquad
    \epsilon\colon 3\to 4.
    \]
    
    \medskip
    \noindent\textbf{Indecomposable projectives.}
    
    \smallskip
    \noindent\emph{The projective $P(4)$.}
    Since $4$ is a sink, one has
    \[
    P(4)=S(4)\colon\quad
    \begin{tikzcd}
    0 \arrow[r] \arrow[rr, bend left] & 0 \arrow[r] \arrow[d] & 0 \arrow[ld] \\
                                      & k &
    \end{tikzcd}.
    \]
    
    \smallskip
    \noindent\emph{The projective $P(3)$.}
    The paths starting at $3$ are
    \[
    e_3 \text{ (ending at $3$)},\qquad \epsilon \text{ (ending at $4$)}.
    \]
    Thus
    \[
    P(3)_1=0,\qquad P(3)_2=0,\qquad P(3)_3\cong k,\qquad P(3)_4\cong k,
    \]
    and the arrow $\epsilon\colon 3\to 4$ induces the identity map. Hence
    \[
    P(3)\colon\quad
    \begin{tikzcd}
    0 \arrow[r] \arrow[rr, bend left] & 0 \arrow[r] \arrow[d] & k \arrow[ld, "1"] \\
                                      & k &
    \end{tikzcd}.
    \]
    
    \smallskip
    \noindent\emph{The projective $P(2)$.}
    The paths starting at $2$ are
    \[
    e_2 \text{ (ending at $2$)},\qquad
    \gamma \text{ (ending at $3$)},\qquad
    \delta,\epsilon\gamma \text{ (ending at $4$)}.
    \]
    Therefore
    \[
    P(2)_1=0,\qquad P(2)_2\cong k,\qquad P(2)_3\cong k,\qquad P(2)_4\cong k^2.
    \]
    Using the ordered basis $\{\delta,\epsilon\gamma\}$ at vertex $4$, we get
    \[
    P(2)\colon\quad
    \begin{tikzcd}
    0 \arrow[r] \arrow[rr, bend left]
    & k \arrow[r, "1"] \arrow[d, "{\begin{psmallmatrix}1\\0\end{psmallmatrix}}"]
    & k \arrow[ld, "{\begin{psmallmatrix}0\\1\end{psmallmatrix}}"] \\
    & k^2 &
    \end{tikzcd}.
    \]
    
    \smallskip
    \noindent\emph{The projective $P(1)$.}
    The paths starting at $1$ are
    \[
    e_1 \text{ (ending at $1$)},\qquad
    \alpha \text{ (ending at $2$)},
    \]
    \[
    \beta,\gamma\alpha \text{ (ending at $3$)},
    \qquad
    \delta\alpha,\epsilon\beta,\epsilon\gamma\alpha \text{ (ending at $4$)}.
    \]
    Hence
    \[
    P(1)_1\cong k,\qquad
    P(1)_2\cong k,\qquad
    P(1)_3\cong k^2,\qquad
    P(1)_4\cong k^3.
    \]
    Choose the ordered bases
    \[
    P(1)_3:\ \{\beta,\gamma\alpha\},
    \qquad
    P(1)_4:\ \{\delta\alpha,\epsilon\beta,\epsilon\gamma\alpha\}.
    \]
    Then
    \[
    P(1)\colon\quad
    \begin{tikzcd}[ampersand replacement=\&]
    k \arrow[r, "1"]
      \arrow[rr, bend left, "{\begin{psmallmatrix}1\\0\end{psmallmatrix}}"]
    \& k \arrow[r, "{\begin{psmallmatrix}0\\1\end{psmallmatrix}}"]
        \arrow[d, "{\begin{psmallmatrix}1\\0\\0\end{psmallmatrix}}"]
    \& k^2 \arrow[ld, "{\begin{psmallmatrix}
    0&0\\
    1&0\\
    0&1
    \end{psmallmatrix}}"] \\
    \& k^3 \&
    \end{tikzcd}.
    \]
    
    \medskip
    \noindent\textbf{Indecomposable injectives.}
    
    \smallskip
    \noindent\emph{The injective $I(1)$.}
    Since $1$ is a source, we have
    \[
    I(1)=S(1)\colon\quad
    \begin{tikzcd}
    k \arrow[r] \arrow[rr, bend left] & 0 \arrow[r] \arrow[d] & 0 \arrow[ld] \\
                                      & 0 &
    \end{tikzcd}.
    \]
    
    \smallskip
    \noindent\emph{The injective $I(2)$.}
    The paths ending at $2$ are
    \[
    e_2 \text{ (starting at $2$)},\qquad \alpha \text{ (starting at $1$)}.
    \]
    Thus
    \[
    I(2)_1\cong k,\qquad I(2)_2\cong k,\qquad I(2)_3=0,\qquad I(2)_4=0,
    \]
    and
    \[
    I(2)\colon\quad
    \begin{tikzcd}
    k \arrow[r, "1"] \arrow[rr, bend left, "0"]
    & k \arrow[r, "0"] \arrow[d, "0"]
    & 0 \arrow[ld] \\
    & 0 &
    \end{tikzcd}.
    \]
    
    \smallskip
    \noindent\emph{The injective $I(3)$.}
    The paths ending at $3$ are
    \[
    e_3 \text{ (starting at $3$)},\qquad
    \gamma \text{ (starting at $2$)},\qquad
    \beta,\gamma\alpha \text{ (starting at $1$)}.
    \]
    Hence
    \[
    I(3)_1\cong k^2,\qquad I(3)_2\cong k,\qquad I(3)_3\cong k,\qquad I(3)_4=0.
    \]
    With respect to the ordered basis $\{\beta,\gamma\alpha\}$ at vertex $1$, we obtain
    \[
    I(3)\colon\quad
    \begin{tikzcd}
    k^2 \arrow[r, "{(0,1)}"] \arrow[rr, bend left, "{(1,0)}"]
    & k \arrow[r, "1"] \arrow[d, "0"]
    & k \arrow[ld, "0"] \\
    & 0 &
    \end{tikzcd}.
    \]
    
    \smallskip
    \noindent\emph{The injective $I(4)$.}
    The paths ending at $4$ are
    \[
    e_4 \text{ (starting at $4$)},\qquad
    \epsilon \text{ (starting at $3$)},
    \]
    \[
    \delta,\epsilon\gamma \text{ (starting at $2$)},
    \qquad
    \delta\alpha,\epsilon\beta,\epsilon\gamma\alpha \text{ (starting at $1$)}.
    \]
    Therefore
    \[
    I(4)_1\cong k^3,\qquad I(4)_2\cong k^2,\qquad I(4)_3\cong k,\qquad I(4)_4\cong k.
    \]
    Choose the ordered bases
    \[
    I(4)_1:\ \{\delta\alpha,\epsilon\beta,\epsilon\gamma\alpha\},
    \qquad
    I(4)_2:\ \{\delta,\epsilon\gamma\}.
    \]
    Then
    \[
    I(4)\colon\quad
    \begin{tikzcd}[ampersand replacement=\&]
    k^3
    \arrow[r, "{\begin{psmallmatrix}
    1&0&0\\
    0&0&1
    \end{psmallmatrix}}"]
    \arrow[rr, bend left, "{(0,1,0)}"]
    \&
    k^2 \arrow[r, "{(0,1)}"] \arrow[d, "{(1,0)}"]
    \&
    k \arrow[ld, "1"] \\
    \& k \&
    \end{tikzcd}.
    \]
    
    \medskip
    \noindent
    This determines all indecomposable projective and injective representations for $Q_1$ and $Q_2$.
    \end{proof}

%=============================================================================
% Problem 15
%=============================================================================
\item Compute a projective resolution of each simple representation $S(i)$ for the quivers $Q_1$ and $Q_2$ from Problem~14.

\begin{proof}
    We compute the (minimal) projective resolutions of the simple representations for the two quivers from Problem~14.
    
    Recall that for an acyclic quiver $Q$, the category $\operatorname{Rep}Q$ is hereditary. Hence every representation has projective dimension at most $1$. In particular, for each simple representation $S(i)$, its minimal projective resolution is obtained from the projective cover
    \[
    P(i)\twoheadrightarrow S(i),
    \]
    whose kernel is $\operatorname{rad} P(i)$.
    
    Moreover, for a path algebra of an acyclic quiver, one has
    \[
    \operatorname{rad} P(i)\cong \bigoplus_{\alpha:i\to j} P(j),
    \]
    where the direct sum runs over all arrows starting at $i$.
    
    \bigskip
    \noindent\textbf{(1) The quiver $Q_1$.}
    \[
    Q_1\colon\quad
    \begin{tikzcd}
    1 & 2 \arrow[l, shift left, "\alpha"] \arrow[l, shift right, "\beta"']
    & 3 \arrow[l, shift left, "\gamma"] \arrow[l, shift right, "\delta"']
    \end{tikzcd}
    \]
    
    \medskip
    \noindent\textbf{Resolution of $S(1)$.}
    Since $1$ is a sink, $P(1)=S(1)$ is already projective. Hence the minimal projective resolution is
    \[
    0 \longrightarrow P(1) \xrightarrow{\;\cong\;} S(1) \longrightarrow 0.
    \]
    
    \medskip
    \noindent\textbf{Resolution of $S(2)$.}
    There are exactly two arrows starting at $2$, namely $\alpha,\beta:2\to 1$. Therefore
    \[
    \operatorname{rad} P(2)\cong P(1)\oplus P(1).
    \]
    Thus we obtain a short exact sequence
    \[
    0 \longrightarrow P(1)\oplus P(1)
    \xrightarrow{\;f\;}
    P(2)
    \xrightarrow{\;\pi\;}
    S(2)
    \longrightarrow 0,
    \]
    where $\pi$ is the canonical projection onto the top of $P(2)$, and $f$ is the morphism whose two components correspond to the two arrows $\alpha$ and $\beta$.
    
    More explicitly, if we denote by
    \[
    \iota_\alpha,\iota_\beta : P(1)\to P(2)
    \]
    the morphisms sending the generator of $P(1)_1\cong k$ to the basis vectors
    \[
    \alpha,\beta \in P(2)_1\cong k^2,
    \]
    then
    \[
    f=(\iota_\alpha,\iota_\beta).
    \]
    Its image is precisely the radical of $P(2)$, so this is the minimal projective resolution of $S(2)$.
    
    \medskip
    \noindent\textbf{Resolution of $S(3)$.}
    There are exactly two arrows starting at $3$, namely $\gamma,\delta:3\to 2$. Hence
    \[
    \operatorname{rad} P(3)\cong P(2)\oplus P(2).
    \]
    Therefore
    \[
    0 \longrightarrow P(2)\oplus P(2)
    \xrightarrow{\;g\;}
    P(3)
    \xrightarrow{\;\rho\;}
    S(3)
    \longrightarrow 0
    \]
    is exact, where $\rho$ is the canonical projection onto the top of $P(3)$, and $g$ is induced by the arrows $\gamma$ and $\delta$.
    
    Explicitly, if
    \[
    \iota_\gamma,\iota_\delta : P(2)\to P(3)
    \]
    are the morphisms sending the generator of $P(2)_2\cong k$ to
    \[
    \gamma,\delta \in P(3)_2\cong k^2,
    \]
    then
    \[
    g=(\iota_\gamma,\iota_\delta).
    \]
    Its image is $\operatorname{rad} P(3)$, so this gives the minimal projective resolution of $S(3)$.
    
    \bigskip
    \noindent\textbf{(2) The quiver $Q_2$.}
    \[
    Q_2\colon\quad
    \begin{tikzcd}
    1 \arrow[r, "\alpha"] \arrow[rr, bend left, "\beta"]
    & 2 \arrow[r, "\gamma"] \arrow[d, "\delta"']
    & 3 \arrow[ld, "\epsilon"] \\
    & 4 &
    \end{tikzcd}
    \]
    
    \medskip
    \noindent\textbf{Resolution of $S(4)$.}
    Since $4$ is a sink, $P(4)=S(4)$ is projective. Hence
    \[
    0 \longrightarrow P(4) \xrightarrow{\;\cong\;} S(4) \longrightarrow 0
    \]
    is the minimal projective resolution of $S(4)$.
    
    \medskip
    \noindent\textbf{Resolution of $S(3)$.}
    There is exactly one arrow starting at $3$, namely $\epsilon:3\to 4$. Thus
    \[
    \operatorname{rad} P(3)\cong P(4),
    \]
    and we get
    \[
    0 \longrightarrow P(4)
    \xrightarrow{\;\iota_\epsilon\;}
    P(3)
    \xrightarrow{\;\pi_3\;}
    S(3)
    \longrightarrow 0,
    \]
    where $\iota_\epsilon$ is the morphism sending the generator of $P(4)_4\cong k$ to
    \[
    \epsilon \in P(3)_4\cong k.
    \]
    This is the minimal projective resolution of $S(3)$.
    
    \medskip
    \noindent\textbf{Resolution of $S(2)$.}
    There are two arrows starting at $2$, namely
    \[
    \gamma:2\to 3,
    \qquad
    \delta:2\to 4.
    \]
    Hence
    \[
    \operatorname{rad} P(2)\cong P(3)\oplus P(4).
    \]
    Therefore
    \[
    0 \longrightarrow P(3)\oplus P(4)
    \xrightarrow{\;h\;}
    P(2)
    \xrightarrow{\;\pi_2\;}
    S(2)
    \longrightarrow 0
    \]
    is exact, where $h=(\iota_\gamma,\iota_\delta)$ and:
    \begin{itemize}
        \item $\iota_\gamma:P(3)\to P(2)$ sends the generator of $P(3)_3\cong k$ to
        \[
        \gamma \in P(2)_3\cong k;
        \]
        \item $\iota_\delta:P(4)\to P(2)$ sends the generator of $P(4)_4\cong k$ to
        \[
        \delta \in P(2)_4\cong k^2.
        \]
    \end{itemize}
    The image of $h$ is exactly the radical of $P(2)$, so this is the minimal projective resolution of $S(2)$.
    
    \medskip
    \noindent\textbf{Resolution of $S(1)$.}
    There are two arrows starting at $1$, namely
    \[
    \alpha:1\to 2,
    \qquad
    \beta:1\to 3.
    \]
    Thus
    \[
    \operatorname{rad} P(1)\cong P(2)\oplus P(3).
    \]
    Hence we have the exact sequence
    \[
    0 \longrightarrow P(2)\oplus P(3)
    \xrightarrow{\;u\;}
    P(1)
    \xrightarrow{\;\pi_1\;}
    S(1)
    \longrightarrow 0,
    \]
    where $u=(\iota_\alpha,\iota_\beta)$ and:
    \begin{itemize}
        \item $\iota_\alpha:P(2)\to P(1)$ sends the generator of $P(2)_2\cong k$ to
        \[
        \alpha \in P(1)_2\cong k;
        \]
        \item $\iota_\beta:P(3)\to P(1)$ sends the generator of $P(3)_3\cong k$ to
        \[
        \beta \in P(1)_3\cong k^2.
        \]
    \end{itemize}
    Its image is $\operatorname{rad} P(1)$, so this yields the minimal projective resolution of $S(1)$.
    
    \bigskip
    \noindent\textbf{Summary.}
    
    For $Q_1$:
    \[
    0\to P(1)\xrightarrow{\cong} S(1)\to 0,
    \]
    \[
    0\to P(1)^2\to P(2)\to S(2)\to 0,
    \]
    \[
    0\to P(2)^2\to P(3)\to S(3)\to 0.
    \]
    
    For $Q_2$:
    \[
    0\to P(4)\xrightarrow{\cong} S(4)\to 0,
    \]
    \[
    0\to P(4)\to P(3)\to S(3)\to 0,
    \]
    \[
    0\to P(3)\oplus P(4)\to P(2)\to S(2)\to 0,
    \]
    \[
    0\to P(2)\oplus P(3)\to P(1)\to S(1)\to 0.
    \]
    
    These are precisely the minimal projective resolutions of the simple representations.
    \end{proof}

\end{enumerate}
\end{document}
